%% =====================================================================
%% main.tex --- The Complement Rule Is Excluded Middle:
%%              Machine-Checked Probability over Intuitionistic Logic
%% Draft for CPP 2027 submission (acmart sigplan).
%%
%% Build: make          (pdflatex + bibtex + 2x pdflatex)
%% =====================================================================

\documentclass[sigplan,10pt,review,anonymous]{acmart}

\settopmatter{printfolios=true,printccs=false,printacmref=false}
\setcopyright{none}
\renewcommand\footnotetextcopyrightpermission[1]{}
\pagestyle{plain}

\usepackage{listingsutf8}
\usepackage{microtype}

% --- Lean 4 listing setup ---------------------------------------------
\lstdefinelanguage{Lean}{
  morekeywords={theorem,lemma,def,noncomputable,structure,instance,
                variable,import,namespace,end,open,scoped,where,
                fun,match,with,if,then,else,let,in,by,do,
                Prop,Type,Sort,deriving,inductive,abbrev,example,
                forall,exists,class,extends},
  sensitive=true,
  morecomment=[s]{/-}{-/},
  morecomment=[l]{--},
  literate=
    {->}{{$\rightarrow$}}{2}
    {→}{{$\rightarrow$}}{1}
    {↔}{{$\leftrightarrow$}}{1}
    {↦}{{$\mapsto$}}{1}
    {⟺}{{$\iff$}}{2}
    {∀}{{$\forall$}}{1}
    {∃}{{$\exists$}}{1}
    {∧}{{$\wedge$}}{1}
    {∨}{{$\vee$}}{1}
    {¬}{{$\neg$}}{1}
    {≤}{{$\le$}}{1}
    {≥}{{$\ge$}}{1}
    {≠}{{$\ne$}}{1}
    {∈}{{$\in$}}{1}
    {∉}{{$\notin$}}{1}
    {⊤}{{$\top$}}{1}
    {⊥}{{$\bot$}}{1}
    {⊔}{{$\sqcup$}}{1}
    {⊓}{{$\sqcap$}}{1}
    {⋁}{{$\bigvee$}}{1}
    {⨆}{{$\bigsqcup$}}{1}
    {∑}{{$\sum$}}{1}
    {∞}{{$\infty$}}{1}
    {ℝ}{{$\mathbb{R}$}}{1}
    {ℕ}{{$\mathbb{N}$}}{1}
    {ℤ}{{$\mathbb{Z}$}}{1}
    {Ω}{{$\Omega$}}{1}
    {δ}{{$\delta$}}{1}
    {μ}{{$\mu$}}{1}
    {ρ}{{$\rho$}}{1}
    {ι}{{$\iota$}}{1}
    {σ}{{$\sigma$}}{1}
    {∂}{{$\partial$}}{1}
    {×}{{$\times$}}{1}
    {⁻¹}{{${}^{-1}$}}{2}
    {ᶜ}{{${}^{\mathsf{c}}$}}{1}
    {ᶜᶜ}{{${}^{\mathsf{cc}}$}}{2}
    {⟨}{{$\langle$}}{1}
    {⟩}{{$\rangle$}}{1}
    {∘}{{$\circ$}}{1}
    {≡}{{$\equiv$}}{1}
    {∅}{{$\emptyset$}}{1}
    {▷}{{$\triangleright$}}{1}
    {⊆}{{$\subseteq$}}{1}
    {≃}{{$\simeq$}}{1}
    {Σ}{{$\Sigma$}}{1}
    {⋃}{{$\bigcup$}}{1}
    {⨅}{{$\bigsqcap$}}{1}
    {∣}{{$\mid$}}{1}
    {β}{{$\beta$}}{1}
    {κ}{{$\kappa$}}{1}
    {Θ}{{$\Theta$}}{1}
    {Γ}{{$\Gamma$}}{1}
    {Δ}{{$\Delta$}}{1}
    {α}{{$\alpha$}}{1}
    {₀}{{$_{0}$}}{1}
    {₁}{{$_{1}$}}{1}
    {₂}{{$_{2}$}}{1}
    {ˢ}{{$^{\mathsf{s}}$}}{1}
    {↑}{{$\uparrow$}}{1}
    {↓}{{$\downarrow$}}{1}
}

\lstset{
  language=Lean,
  inputencoding=utf8/latin1,
  extendedchars=true,
  basicstyle=\ttfamily\small,
  keywordstyle=\bfseries,
  commentstyle=\itshape,
  showstringspaces=false,
  breaklines=true,
  xleftmargin=1em,
  aboveskip=0.6em,
  belowskip=0.6em,
}

\acmConference[CPP '27]{Certified Programs and Proofs}{January 2027}{(submitted)}
\acmYear{2027}
\acmISBN{}
\acmDOI{}
\citestyle{acmauthoryear}

\theoremstyle{plain}
\newtheorem{theorem}{Theorem}
\newtheorem{lemma}[theorem]{Lemma}
\newtheorem{proposition}[theorem]{Proposition}
\newtheorem{corollary}[theorem]{Corollary}
\theoremstyle{definition}
\newtheorem{definition}[theorem]{Definition}
\newtheorem{example}[theorem]{Example}
\theoremstyle{remark}
\newtheorem{remark}[theorem]{Remark}

% amsthm wraps the theorem-name argument in \mbox by default, which makes
% it unbreakable and lets long \texttt{} identifiers overflow the narrow
% two-column width even with \allowbreak inserted.  Drop the \mbox so the
% note breaks like ordinary text.
\renewcommand{\thmnote}[1]{\normalfont #1}

\emergencystretch=3em

% Shorthands
\newcommand{\ENNReal}{\ensuremath{[0,\infty]}}
\newcommand{\slack}{\mathrm{slack}}
\newcommand{\Bel}{\mathrm{Bel}}
\newcommand{\Pl}{\mathrm{Pl}}

% Order-theoretic notation (fixed in Section 4.1): principal lower set,
% strict principal lower set, principal upper set.
\newcommand{\pdown}{{\downarrow}}
\newcommand{\pdowns}{\mathop{\downarrow\!\!\downarrow}}
\newcommand{\pup}{{\uparrow}}

% Link from a main-text result to its Lean statement in the correspondence
% appendix (app-lean.tex).  Renders as a small clickable superscript marker.
\newcommand{\leanref}[1]{\hyperref[lean:#1]{\textsuperscript{\textsf{\tiny Lean}}}}

\begin{document}

\title{Extending Intuitionistic Logic to a Probability Calculus (in Lean 4)}

\author{Anonymous Author(s)}
\affiliation{\institution{Anonymous Institution}\city{}\country{}}

\begin{abstract}
Kolmogorov's probability fixes the event algebra to be Boolean, so
every event is decidable.
This is a poor fit for the semidecidable properties, termination,
reachability, that verification cares about. 
Following the Cox--Jaynes interpretation of probability as extended logic,
we define a construct for extending a logic (treated as a parameter in 
our approach) to a calculus of probability, 
and specify criteria for when this works: (1) when the resulting calculus 
can reconstruct the original logic in the limiting case of certain 
($0-$ or $1-$ probability) events, and (2) when the construction, evaluated 
at classical logic, agrees with classical probability theory.
Instantiating the construction at intuitionistic logic gives valuations
on frames, complete Heyting algebras of open sets, the algebras of
verifiable propositions. In this setting,
we define slack, $1 - v(A) - v(\neg A)$, the classical complement
rule's deficit and boundary mass in the measure model.  Slack
pinpoints exactly when Dempster and Bayesian conditioning disagree.
A separate obstruction, a five-element counterexample, is why Cox's
sum rule cannot be derived and must instead be posited.
We show that a classical, slack-free readout of a semidecidable property
is uncomputable: computing it uniformly is exactly as hard as deciding
the halting problem.
The whole development is machine-checked in Lean~4 over \textsf{mathlib},
in roughly $5{,}000$ lines with no \texttt{sorry}.
\end{abstract}

\begin{CCSXML}
<ccs2012>
<concept>
<concept_id>10003752.10003790.10003792</concept_id>
<concept_desc>Theory of computation~Proof theory</concept_desc>
<concept_significance>500</concept_significance>
</concept>
<concept>
<concept_id>10003752.10003790.10011740</concept_id>
<concept_desc>Theory of computation~Type theory</concept_desc>
<concept_significance>300</concept_significance>
</concept>
<concept>
<concept_id>10002950.10003648</concept_id>
<concept_desc>Mathematics of computing~Probability and statistics</concept_desc>
<concept_significance>500</concept_significance>
</concept>
</ccs2012>
\end{CCSXML}
\ccsdesc[500]{Theory of computation~Proof theory}
\ccsdesc[500]{Mathematics of computing~Probability and statistics}

\keywords{intuitionistic logic, probability, Cox's theorem, belief
functions, locales, Lean, mathlib, formalized mathematics}

\maketitle

\section{Introduction}
\label{sec:intro}


Kolmogorov's probability~\cite{kolmogorov1933} assigns numerical values representing 
uncertainty to events. The event logic is classical, forming a Boolean $\sigma$-algebra.
The assignments satisfy non-negativity, normalization, and $\sigma$-additivity. This commits 
the framework to the axioms of classical logic, including the law of excluded middle. Many 
properties of programs (e.g. "does this computation halt?", "does this process reach a given 
state?") are only \emph{semidecidable}: a procedure confirms the property when it holds, and 
no procedure, in general, confirms its failure.
The mismatch that results is not a contradiction. It is a failure of discrimination. 
A probability calculus built on a particular logic inherits that logic's identifications, so it 
cannot discriminate events the logic proves equal. 

In the case of classical logic, this identification is the \emph{law of excluded middle}:
$\mathit{halts} \sqcup \neg\mathit{halts} = \top$, so the Boolean event algebra contains 
$\neg\mathit{halts}$. The complement rule $P(A^{c}) = 1 - P(A)$ then forces 
the remainder $1 - P(\mathit{halts})$ onto $\neg\mathit{halts}$, an event no finite 
observation can confirm. The events ``not yet confirmed'' and ``refuted'' are identified.
For a program whose non-termination is independent of ZFC~\cite{yedidia-aaronson2016},
the classical value $P(\mathit{halts}) = 0$ is correct, yet neither running the program
nor any ZFC proof ever establishes it. Uniformly over programs, computing these sharp
values is exactly as hard as the halting problem. 

Another failure of discrimination appears in \emph{linear logic}, whose
tensor $\otimes$ distinguishes $a$ from $a \otimes a$, one use of a
resource from two. Modularity with $\otimes$ in place of $\sqcap$,
$v(a) + v(b) = v(a \sqcup b) + v(a \otimes b)$, inherits idempotence of
disjunction, $a \sqcup a = a$, and at $b = a$ forces
$v(a \otimes a) = v(a)$. So, no
assignment satisfying this axiom distinguishes one use of $a$ from two.

Probability is, as the Cox--Jaynes tradition has it, 
\emph{extended logic}~\cite{cox1946,jaynes2003}, and different logics can be extended 
through the same construction~\cite{paris2001,williams2012}. The logic of events, in its algebraic
form the event algebra, is a \emph{parameter} of the construction, and
the choice of event algebra determines the probability calculus.
An event algebra consists of a set of events together with operations 
interpreting the logic's connectives and structural rules, while a probability 
calculus is differentiated from an arbitrary grading of events by satisfying 
certain axioms. We specify how to construct such extensions, and
propose two criteria for when the extension of a logic is a probability
calculus:

\begin{itemize}
  \item[(1)] the valuations taking only the values $0$ and 
$1$ (the two degrees of complete certainty) must reconstruct the underlying event logic.
This is what we take ``extended logic'' to mean.
  \item[(2)] the construction evaluated at 
classical logic must return the classical theory, ensuring it is a genuine generalization 
of classical probability.
\end{itemize}

We propose two fixes for the identifications above. Both extend a logic
to a probability calculus, and on finite frames the second reduces to
the first. 
The first fix we propose is one that accepts the standard axioms imposed on a grading 
to make it a probability calculus, while choosing an event algebra that 
makes the required distinctions. 
The extension of the event algebra to a probability calculus is then given by 
the \emph{valuations}, i.e. the maps from events to numbers satisfying monotonicity, 
normalization, and modularity~\cite{weatherson2003,paris2001}. The resulting extension does not guarantee
defining a probability calculus, and we specify the conditions under which it
does: the $\{0,1\}$-valued valuations are the points of the event algebra,
and they reconstruct its order when it is distributive. 
Intuitionistic logic, in
particular, is a case for which this approach yields an extension, and we
work it out in detail. Its event algebras are the Heyting algebras, and we 
work with frames (the complete Heyting algebras). The open sets of any topological space
form a frame. On a frame the pseudocomplement $\neg a$ need not satisfy
$a \sqcup \neg a = \top$, so the complement rule weakens to
$v(a) + v(\neg a) \le 1$. The gap
$\slack(a) = 1 - v(a) - v(\neg a)$ is the mass of the part of $a$
neither confirmed nor refuted. For the valuation that a measure $\mu$
induces on the open sets of a space, the slack of an open set is the
$\mu$-measure of its boundary.

The second fix also applies to logics whose connectives are not lattice
operations, such as linear logic. It generalizes the axioms of the
probability calculus in addition to varying the event algebra. It
applies to any logic presented by a set of \emph{certain states}, its
$\{0,1\}$-valued assignments. The calculus is the set of finite mixtures
of certain states, the construction of de Finetti that generalized
probabilism applies to non-classical logics~\cite{paris2001,williams2012}.
A mixture that takes only the values $0$ and $1$ is
one of the certain states it mixes, because every certain state with
positive weight must agree with it at every event. So the
$\{0,1\}$-valued elements of the calculus are exactly the certain
states, which present the logic, and criterion (1) holds. 
For finitely many events, the axioms of
the calculus are exactly the affine laws that hold on every certain
state. When the certain states respect $\wedge$ and $\vee$, Paris shows that these
laws are monotonicity, normalization, and modularity~\cite{paris2001}, so
on finite frames and finite Boolean algebras the calculus is the set of
valuations. 

For linear logic, the event algebra is a quantale, a complete lattice
with an associative tensor $\otimes$ that need not be idempotent. The
first fix fails there because of the probability axioms, not the event
algebra: the modularity axiom identifies $a$ with $a \otimes a$.
Probabilistic coherence spaces~\cite{danos-ehrhard2011} avoid this by
assigning weights to the points of each formula's web instead of to
formulas, with $a \otimes a$ interpreted on a product web. Their
elements are not valuations in the sense of the modularity axiom. For
every multiplicative--additive formula over finite data types, the
calculus of the second fix is contained in the probabilistic coherence
space of the formula, and the certain elements of that space are cliques
of Girard's coherence space~\cite{girard1987}. 

The calculus and the probabilistic coherence space coincide at
first-order linear logic. At higher order they can
differ: for a program with two functional arguments the calculus is the
set of randomized sequential programs, a strict subset of the
probabilistic coherence space (Section~\ref{sec:second}). The probabilistic
coherence space, the standard model of probabilistic programs at this
type, therefore admits behaviours that no sequential program flipping coins can
produce in a linear probabilistic language for this type. Reasoning in
it can certify only weaker guarantees than the programs satisfy. The gap
has the same numerical form as the gap between classical and
non-classical correlations in quantum foundations~\cite{cabello2014}. At
$!\mathrm{Bool}$ the calculus is also a strict subset.

Quantum logic, whose event algebras are orthomodular
lattices~\cite{birkhoff-vonneumann1936}, is a logic to which neither fix
applies. Both rely on $\{0,1\}$-valued assignments: the first on
$\{0,1\}$-valued valuations, which need not exist on non-distributive
lattices and do not exist on $M_3$, and the second relies on having certain states. A certain
state of quantum logic would assign $1$ to exactly one ray of every
orthogonal basis, and by the Kochen--Specker
theorem~\cite{kochen-specker1967} there is none in dimension at least
$3$. We prove this for an $18$-ray configuration in
$\mathbb{R}^4$~\cite{cabello1996}, so the calculus of the second fix is
empty there. Quantum states, which by Gleason's
theorem~\cite{gleason1957} are given by density matrices, nevertheless
exist, and they fail criterion (1).


\paragraph{Contributions.}

\begin{enumerate}
\item \textbf{Characterizations} (Sections~\ref{sec:valuation},
  \ref{sec:hinge} and~\ref{sec:ds}).  The complement rule holds for
  every valuation iff the frame satisfies excluded middle, Dempster's
  rule agrees with Bayesian conditioning iff the evidence has zero
  slack, and slack splits exactly into a double-negation part and a
  De Morgan part.

\item \textbf{Structure theorems} (Sections~\ref{sec:representation}
  and~\ref{sec:ds}).  Every valuation on the lower sets of a poset
  splits into an atomic part carried by points and a diffuse remainder, and on its Boolean core it
  is a belief function.

\item \textbf{Cox's theorem} (Section~\ref{sec:cox}).  The sum-rule
  half admits no derivation over a frame, even with weak excluded
  middle, since on a five-element frame a disjunction's value is not a
  function of its disjuncts' values.

\item \textbf{Grounding models} (Section~\ref{sec:bridges}).  For a
  classical measure read on the opens, slack is boundary mass, and a
  classical posterior overstates refutation by exactly the conditional
  boundary mass $\mu(\partial A \cap B)/\mu(B)$.

\item \textbf{Products} (Section~\ref{sec:monad}).  Product valuations
  are unique on $\mathrm{Low}(\mathbb{N} \times \mathbb{N})$, and the
  countably-presented valuations compose associatively and
  commutatively.

\item \textbf{Beyond frames} (Section~\ref{sec:lattice}).  Lattice
  valuations are the normalized monotone homomorphisms out of a
  valuation monoid in which every element is a sum over a chain, and
  no modular grading of linear logic's $\otimes$ distinguishes $a$
  from $a \otimes a$.

\item \textbf{Calculi from certain states} (Section~\ref{sec:recipe}).
  The completeness theorem of generalized probabilism has a proof
  without choice for rational points, and on arbitrary sets of events
  the functions satisfying its laws are the barycenters of measures
  concentrated on the certain states.

\item \textbf{Linear logic} (Sections~\ref{sec:recipe-linear}
  and~\ref{sec:recipe-mall}).  The calculus of a
  multiplicative--additive formula is contained in its probabilistic
  coherence space, whose certain elements are cliques, with equality
  at implications between data types and strict inclusion at
  $!\mathrm{Bool}$.

\item \textbf{Second-order linear types} (Section~\ref{sec:second}).
  For a program with two functional arguments over data types with
  at least two elements, the calculus is the
  stable-set polytope of a graph whose clique-constrained polytope is
  the probabilistic coherence space, and the resulting strict inclusion
  transfers to every formula in which this type occurs positively.

\item \textbf{Quantum event logic} (Section~\ref{sec:quantum}).  On
  Cabello's Kochen--Specker set the calculus is empty, and with
  divergence allowed every mixture of certain states has total weight
  at most $4$, against $\tfrac92$ for every quantum state.

\item \textbf{A first formal treatment.}  Every result above is machine-checked in
  Lean~4 over \textsf{mathlib}, except cited theorems the text marks as
  not formalized, which to our knowledge is the first mechanization of
  non-additive probability over intuitionistic logic.
\end{enumerate}

\paragraph{Mechanization.}
The development is carried out in Lean~4~\cite{moura-ullrich2021} over
\textsf{mathlib}~\cite{mathlib2020}, in roughly $12{,}000$ lines with no
\texttt{sorry}, and is included in full as anonymous supplementary
material.\footnote{Claude (Anthropic) was used to assist with
drafting and debugging the Lean formalization.} Two points are specific
to a proof assistant.  We separate \emph{object-level} classicality,
excluded middle in the frame $\Omega$, from the \emph{meta-level}
classicality \textsf{mathlib} inherits through \texttt{Classical.choice},
identify where the meta-level dependence is essential
(Table~\ref{tab:classical}), and give proofs without choice of the
finite results that admit them.  
The product-rule half of Cox's theorem needs an
Acz\'el/H\"older regraduation absent from \textsf{mathlib}, which we
build from scratch.

\section{Valuations on Frames}
\label{sec:valuation}

\subsection{The Event Logic}

\begin{definition}[frame]
\label{def:frame}
A \emph{frame} is a complete lattice $\Omega$ in which finite meets
distribute over arbitrary joins,
\[
  a \sqcap \Bigl(\bigsqcup_{i \in I} b_i\Bigr)
  \;=\; \bigsqcup_{i \in I} (a \sqcap b_i)
\]
for every $a \in \Omega$ and every family $\{b_i\}_{i \in I}
\subseteq \Omega$.  Equivalently, it is a complete Heyting algebra.
By $x \Rightarrow y$ we denote the Heyting implication in a frame,
and $\neg a := a \Rightarrow \bot$ is the pseudocomplement of
$a \in \Omega$.  Also, $a^{c}$ denotes a complement of $a$, if it
exists, i.e.\ an element such that $a \sqcap a^{c} = \bot$ and
$a \sqcup a^{c} = \top$.
\end{definition}

\noindent
Frames are the algebraic semantics of
intuitionistic logic and the point-free abstraction of topological
spaces.  The same objects with morphisms reversed form the category
of \emph{locales}, so the distinction is invisible at the level of
a single algebra.  We say \emph{frame} throughout and reserve
\emph{locale} for the spatial and categorical statements (points of
a locale, classifying locales) where that word is the standard one.  The motivating example throughout is
$\Omega = \mathcal{O}(X)$, the opens of a space, with
$\sqcap, \sqcup$ as intersection and union and the pseudocomplement
$\neg a$ of an open given by the \emph{interior} of its set
complement.  In \textsf{mathlib} this is the typeclass
\texttt{Order.Frame}, and the Heyting operations come with the key
asymmetry already in place: $a \sqcap \neg a = \bot$ holds in
any Heyting algebra (\texttt{inf\_compl\_eq\_bot}), while
$a \sqcup \neg a = \top$ is available only in Boolean
algebras.  An element with $a \sqcup \neg a = \top$ is
\emph{complemented}: it satisfies its own instance of excluded
middle.

\subsection{The Valuation Structure}

The central object assigns a degree of certainty to each element of
the frame, with inclusion--exclusion as the only interaction axiom:

\begin{definition}[valuation]
\label{def:valuation}
A \emph{valuation} on a frame $\Omega$ is a \emph{monotone} function
$v : \Omega \to \ENNReal$ with
$v(\bot) = 0$, $v(\top) = 1$ (\emph{normalization}), and, for all
$a, b \in \Omega$, \emph{modularity}:
\[
  v(a) + v(b) \;=\; v(a \sqcup b) + v(a \sqcap b).
\]
\end{definition}

\begin{lstlisting}
structure Valuation (Ω : Type*) [Order.Frame Ω] where
  toFun : Ω → ℝ≥0∞
  map_bot' : toFun ⊥ = 0
  map_top' : toFun ⊤ = 1
  mono'    : Monotone toFun
  modular' : ∀ a b, toFun a + toFun b
             = toFun (a ⊔ b) + toFun (a ⊓ b)
\end{lstlisting}

This axiom set is not ours: Weatherson's axioms for intuitionistic
probability are exactly these~\cite{weatherson2003}, and the
valuations of constructive measure
theory~\cite{vickers2008,coquand-spitters2009} are their
Scott-continuous instances (Section~\ref{sec:related}).

Values live in the extended non-negative reals $\ENNReal$, and
normalization confines them to $[0,1]$\leanref{valuation-le-one}.  The
conventions this choice carries, unconditional summation and truncated
subtraction, are an artifact of the mechanization and are recorded in
Appendix~\ref{app:ennreal}.

\subsection{Slack}

Modularity applied to the disjoint pair $(a, \neg a)$
collapses to
$v(a) + v(\neg a) = v(a \sqcup \neg a)$\leanref{valuation-add-compl-eq-sup}.
So, the join need not be $\top$:

\begin{theorem}[the complement bound]
\label{thm:add-compl-le-one}
For every valuation $v$ on a frame $\Omega$ and every $a \in \Omega$,
$v(a) + v(\neg a) \le 1$.
\end{theorem}

\begin{lstlisting}
theorem Valuation.add_compl_le_one (v : Valuation Ω) (a : Ω) :
    v a + v aᶜ ≤ 1
\end{lstlisting}

\noindent
This is a \emph{derived} Dempster--Shafer sub-additivity inequality.
The deficit
\[
  \slack(a) \;=\; 1 - \bigl(v(a) + v(\neg a)\bigr)
  \;=\; 1 - v(a \sqcup \neg a)
\]
is the valuation of the undecided region.  It decomposes canonically.
The pair $(\neg\neg a, \neg a)$, double negation
against negation, is always disjoint. However, it is a partition only under
\emph{weak} excluded middle
($\neg a \sqcup \neg\neg a = \top$).
Applying modularity to this disjoint pair splits
the slack into two logically independent gaps:

\begin{definition}[regular element]
An element $a \in \Omega$ is \emph{regular} (or $\neg\neg$-stable) if
$a = \neg\neg a$.
\end{definition}

\begin{theorem}[slack decomposition\leanref{valuation-slack-eq-dngap-add-demorgangap}]
For every valuation $v$ on a frame $\Omega$ and every $a \in \Omega$,
$\slack(a) = \mathrm{dnGap}(a) + \mathrm{deMorganGap}(a)$, where
$\mathrm{dnGap}(a) = v(\neg\neg a) - v(a)$ is $0$ if $a$ is
$\neg\neg$-stable (regular), and
$\mathrm{deMorganGap}(a) = 1 - v(\neg\neg a \sqcup \neg a)$
is $0$ if weak excluded middle holds at $a$.
\end{theorem}

\noindent
The single Dempster--Shafer ``ignorance'' number is thus resolved
into a regularity defect and a De~Morgan defect: the first measures
failure of double-negation stability, the second corresponds to a
named intermediate logic (weak excluded middle).  
Constructively, the contrapositive of a universal
hypothesis returns only its double-negated form, so for open
sets, this specification of probability does not fall victim
to Hempel's
raven paradox~\cite{hempel1945}: evidence of the white
sock indoors confirms at most $\neg\neg a$, not $a$ itself, and
$\mathrm{dnGap}$ is exactly the gap between confirming
$\neg\neg a$ and confirming $a$ itself.

Because $\mathrm{dnGap}$ and $\mathrm{deMorganGap}$ are themselves
defined from $v$, the apparent three unknowns $v(a)$, $v(\neg a)$, and
$v(\neg\neg a)$ are not independent: they are part of a three-constraint 
system, yielding the following result:

\begin{theorem}[slack in closed form\leanref{valuation-slack-eq-sub-add-sub}]
\label{thm:slack-explicit}
For every valuation $v$ on a frame $\Omega$ and every $a \in \Omega$,
\[
  \slack(a) \;=\; \bigl(v(\neg\neg a) - v(a)\bigr)
  \;+\; \bigl(1 - v(\neg\neg a) - v(\neg a)\bigr),
\]
where $v(a) \le v(\neg\neg a)$.
\end{theorem}

\noindent
The inequality follows from monotonicity applied to
$a \le \neg\neg a$, a fact of every Heyting algebra, and shows the
first summand is exactly the difference, not the floor
produced by $\ENNReal$'s truncated subtraction.  The split into a
regularity defect and a De~Morgan defect follows from monotonicity
and modularity alone, with no further axiom about ignorance required.
This contrasts with the usual Dempster--Shafer picture, where the
mass assigned to non-singleton, ambiguous sets is an extra primitive.

Theorem~\ref{thm:slack-explicit} expresses $\slack$ pointwise
in terms of $v$.  The converse, whether $\slack$ in turn determines
$v$, is more delicate.  One constraint is automatic:

\begin{theorem}[slack is submodular\leanref{valuation-slack-submodular}]
\label{thm:slack-submodular}
For every valuation $v$ on a frame $\Omega$ and every $a, b \in \Omega$,
\[
  \slack(a \sqcup b) + \slack(a \sqcap b) \;\le\; \slack(a) + \slack(b).
\]
\end{theorem}

\noindent
The proof runs $v$'s own modularity twice, once at $(a, b)$ and once
at $(\neg a, \neg b)$, and combines the two results with the De~Morgan
laws: the equality $\neg(a \sqcup b) = \neg a \sqcap \neg b$, which
holds in every Heyting algebra, and the inequality
$\neg a \sqcup \neg b \le \neg(a \sqcap b)$, monotone under $v$ but
not in general an equality.  Any function proposed as the slack of
some valuation must therefore be submodular.

\begin{remark}[non-recovery\leanref{exists-ne-valuation-same-slack}]
\label{rem:no-recovery}
Submodularity is necessary, not sufficient:
Section~\ref{sec:no-recovery} exhibits two valuations on the same
frame with identical slack functions.
\end{remark}

\begin{remark}[slack under weak excluded middle\leanref{valuation-slack-eq-dngap-of-isdemorgan}]
The two gaps, $\mathrm{dnGap}$ and $\mathrm{deMorganGap}$, computed 
at a single element, are logically independent.
However, if
weak excluded middle is imposed on the frame for every $a$
($\neg\neg a \sqcup \neg a = \top$\leanref{isdemorgan}), then
$\mathrm{deMorganGap}$ vanishes identically and
$\slack(a) = \mathrm{dnGap}(a)$ for every valuation and every $a$.
This sits strictly between plain intuitionistic logic, where both
gaps can be positive, and full excluded middle, where both vanish\leanref{classical-slack-zero}.
A concrete instance appears in Section~\ref{sec:bridges}.
\end{remark}

\subsection{What Survives without Excluded Middle}

Three classical pillars survive intact, and one fails in a precisely
delimited way.  \emph{Disjoint additivity} is unconditional:
$a \sqcap b = \bot$ implies $v(a \sqcup b) = v(a) + v(b)$.
\emph{Conditioning} is native: for $v(b) \neq 0$ the posterior
$v(a \mid b) = v(a \sqcap b)/v(b)$ is again a valuation (the proof is
frame distributivity), and the product rule
$v(a \mid b) \cdot v(b) = v(a \sqcap b)$ and Bayes symmetry
(and therefore Bayesian updating)
$v(a \mid b)\,v(b) = v(b \mid a)\,v(a)$ hold\leanref{condval-mul}.  \emph{Total
probability} holds over any genuine partition, meaning
$a \sqcap a' = \bot$ and $a \sqcup a' = \top$:
$v(b) = v(a \sqcap b) + v(a' \sqcap b)$.  What fails is total
probability over $\{a, \neg a\}$.  That family need not join to
$\top$, and the conditional masses fall short of $v(b)$ by a
conditional slack\leanref{cond-add-compl-le}.

\paragraph{The classical fragment.}
We characterize the elements of $\Omega$ on which the classical
complement rule holds.  The $\neg\neg$-stable (\emph{regular})
elements, which form a Boolean algebra (by Glivenko's theorem), are a
natural candidate, but they do not suffice. Modularity constrains
the frame join $\sqcup$, whereas the Boolean join of regular
elements is the regularized $\neg\neg(a \sqcup b)$, for which
only the inequality
$v(a) + v(b) \le v(\neg\neg(a \sqcup b)) + v(a \sqcap b)$ survives
(i.e. the 2-monotonicity of Section~\ref{sec:ds}).  A regular but uncomplemented
element, which is an open ray $(-\infty,0)$ in the opens of $\mathbb{R}$,
under an atomic measure at $0$, still carries slack.  The correct
classical fragment is the \emph{complemented} elements, where
$v(a) + v(a^{c}) = 1$ holds exactly\leanref{add-compl-eq-one-of-complemented}.  On a connected
locale that fragment can be as small as $\{\bot, \top\}$.  In the
Boolean limit, every element is complemented, the slack vanishes
identically, and classical Kolmogorov axioms are recovered\leanref{classical-additivity}.

\paragraph{Mixtures.}
Valuations are closed under finite convex combination\leanref{valuation-mix}, with every axiom inherited termwise,
because modularity is linear.  This innocuous fact becomes a
characterization in Section~\ref{sec:representation}.

\section{Classicality in Theory and Metatheory}
\label{sec:hinge}

Classicality enters this work twice: in the metatheory, the logic
of the proof assistant itself, and in the object theory, the frame's
own logic.

\paragraph{Lean's meta-theory.}
Lean's foundations include
\texttt{Classical.choice}, and \textsf{mathlib} is built on it
throughout.  For example, the order relation on 
$\mathbb{R}$ and $\ENNReal$ has no
computable \texttt{Decidable} instance in general, so any library
lemma that case-splits on $x \le y$ for reals resolves that split via
\texttt{Classical.propDecidable}. A
formalization that never invokes excluded middle at the level of
frame elements $a \in \Omega$ can therefore still be classical 
through the arithmetic it borrows from the library.  We track
this separately from object-level classicality throughout the paper.
Theorem~\ref{thm:decomposition} of Section~\ref{sec:representation}
rests only on this ambient baseline and contains no explicit case
split. The Scott-continuity recovery theorem
(Theorem~\ref{thm:recovery})
invokes Lean's \texttt{classical} tactic directly, and the
explicit Cox model of Section~\ref{sec:cox} goes the other
way, replacing the default with an explicit computable
\texttt{Decidable} instance to keep that one instance genuinely
non-classical.

\paragraph{Object-level classicality: the R3 axiom.}
Van Horn's repaired Cox theorem~\cite{vanhorn2003} axiomatizes
plausibility over classical propositions.  Its negation axiom, called
\emph{R3}, posits that the plausibility of $\neg A$ is a fixed
non-increasing function $S$ of the plausibility of $A$.  Consistency
forces the involution $S \circ S = \mathrm{id}$, and therefore the
complement rule $p \mapsto 1 - p$.  We show that the 
involutive negation \emph{is}
double-negation elimination, so R3 is where classical logic enters
every known Cox-style derivation.  To do this, we first 
isolate the exact content of
R3's conclusion as a predicate on valuations, then prove equivalence
to excluded middle.

That predicate is \emph{classical negation}: a valuation $v$ has
classical negation if $v(a) + v(\neg a) = 1$ for every $a \in \Omega$.

\begin{theorem}[the R3 characterization\leanref{hinge}]
\label{thm:hinge}
For every frame $\Omega$, every valuation on $\Omega$ has classical
negation \emph{iff} $a \sqcup \neg a = \top$ for every
$a \in \Omega$.
\end{theorem}

The forward direction is a
calculation.  The content is the converse: wherever excluded
middle fails, a valuation with genuine slack must be
\emph{constructed}.  Suppose $a \sqcup \neg a \neq \top$.

\begin{definition}[ideal, prime ideal]
An \emph{ideal} $J \subseteq \Omega$ is downward closed and closed
under finite joins.  $J$ is \emph{prime} if $\top \notin J$ and
$a \sqcap b \in J$ implies $a \in J$ or $b \in J$.
\end{definition}

The prime ideal separation theorem
(\textsf{mathlib}'s \texttt{DistribLattice.\allowbreak prime\_ideal\_of\_%
\allowbreak disjoint\_\allowbreak filter\_ideal}) yields a prime ideal $J$ containing
$a \sqcup \neg a$ but not $\top$.  The indicator of $J$'s
complement,
\[
  v(x) = \begin{cases} 0 & x \in J \\ 1 & x \notin J, \end{cases}
\]
is a valuation.  Modularity in the only nontrivial case ($x, y
\notin J$) is \emph{exactly} primeness, $x \sqcap y \in J
\Rightarrow x \in J \vee y \in J$.  And this valuation assigns
$v(a) + v(\neg a) = 0 \neq 1$.

\paragraph{Sharp valuations are the points of the locale.}
The construction above is one half of an exact
correspondence.  Call $v$ \emph{sharp} if it is $\{0,1\}$-valued, a
state of full certainty.  Sharpness characterizes exactly the
prime-ideal indicators:

\begin{theorem}[sharp valuations are points\leanref{sharp-iff-point}]
\label{thm:sharp-iff-point}
For every valuation $v$ on a frame $\Omega$, $v$ is sharp iff $v$ is
the complement-indicator of a prime ideal.
\end{theorem}

\noindent
The correspondence is given by two mutually inverse maps: $v \mapsto Z(v) =
\{x \in \Omega : v(x) = 0\}$, which computes the zero-set of a sharp valuation, 
and the map taking
$J$ to its complement-indicator (the valuation of
Theorem~\ref{thm:hinge}'s construction).  $Z(v)$ is prime exactly
because $v$ is sharp, and the complement-indicator of $Z(v)$
recovers $v$ itself.  In locale-theoretic terms, sharp valuations
recover the \emph{points} of the locale.  A prime ideal is
\emph{completely prime} when it is also closed under arbitrary
(not just finite) suprema of its members. 
The finitely prime points 
of Theorem~\ref{thm:sharp-iff-point} sharpen to
the completely prime points exactly where Scott continuity is added,
recovering points of the locale in the sense used
by~\cite{vickers2008,coquand-spitters2009}:

\begin{theorem}[Scott-continuous sharp valuations\leanref{sharp-scott-iff-completelyprimepoint}]
For every valuation $v$ on a frame $\Omega$, $v$ is sharp and
Scott-continuous iff $v$ is the complement-indicator of a completely
prime ideal.
\end{theorem}

\noindent  The characterization and the two
correspondences together locate the object-level classicality
precisely: whenever
excluded middle fails at some $a$, the failure occurs at a
point of the locale.

\section{Valuation Representation}
\label{sec:representation}

Every valuation on a finite frame is a probability mass function on
points. On infinite frames, this can fail, and Scott continuity
marks the boundary between where it still holds and where
it does not.

\subsection{Finite Frames Representation}

\begin{theorem}[Birkhoff duality~\cite{birkhoff1937}]
Every finite distributive lattice is isomorphic to the lattice of
lower sets of a finite poset: its sup-irreducible elements form a
poset $P$ whose lower sets reconstruct the lattice, and every finite
poset arises this way from its own lattice of lower sets.
\end{theorem}

\noindent
Let $P$ be the finite poset of
sup-irreducible elements of a finite frame.
Then, by Birkhoff duality
(\texttt{OrderIso.lowerSetSupIrred}),
every finite frame is the lattice of \emph{lower sets} of $P$.
We fix the following notation, used throughout the rest of the paper.
A \emph{lower set} $L \subseteq P$ satisfies $p \in L,\, q \le p
\Rightarrow q \in L$, and an \emph{upper set} is defined dually.
Lower sets are always denoted $L$, and upper sets $U$ (opens of a
topological space are still $U, V, W$, as in traditional usage).  The
\emph{principal lower set} at $p$ is $\pdown p = \{q : q \le p\}$, the
\emph{strict principal lower set} at $p$ is
$\pdowns p = \{q : q < p\}$, and the \emph{principal upper set} at
$p$ is $\pup p = \{q : q \ge p\}$.
We prove the finite representation theorem below on the lower sets of
$P$.
Transporting
it across the isomorphism above then gives the theorem for every
finite frame, not just lower sets\leanref{valuation-exists-pmf-presentation}.
Let us define the \emph{mass} of a
point as the discrete derivative of $v$\leanref{valuation-mass}:
\[
  \mathrm{mass}(p) \;=\; v(\pdown p) - v(\pdowns p).
\]

The masses recover $v$ and sum to $1$:

\begin{theorem}[finite representation\leanref{valuation-eq-sum-mass}]
\label{thm:finite-rep}
For a finite $P$, any valuation $v$ on $\mathrm{Low}(P)$, and any
$L \in \mathrm{Low}(P)$,
$v(L) = \sum_{p \in L} \mathrm{mass}(p)$, and
$\sum_{p} \mathrm{mass}(p) = 1$.
\end{theorem}

The proof works by induction on $|L|$: repeatedly remove a maximal
point $p$ from $L$, and use modularity applied to
$(L \setminus \pup p,\; \pdown p)$ to show each removal
changes $v$ by exactly $\mathrm{mass}(p)$.
The mass function is a classical probability mass function on
$P$\leanref{valuation-topmf}.  And since lower sets are
exactly the opens of the Alexandrov topology, where interior is the
identity, the theorem also reads $v = \mu \circ \mathrm{int}$, i.e.
every valuation on a
finite frame is the classical measure $\mu$ applied after
taking the interior.
In the finite case, this is the converse of
the measure-theoretic bridge of Section~\ref{sec:bridges}.
The intuitionistic mass also agrees with the classical notion of the
mass of a point: reading $v = \mu \circ \mathrm{int}$,
\[
  \mathrm{mass}(p) \;=\; \mu(\pdown p) - \mu(\pdowns p)
  \;=\; \mu(\{p\}),
\]
since $\pdown p$ is the disjoint union of $\pdowns p$ and $\{p\}$ and
$\mu$ is defined on all subsets of $P$.

Theorem~\ref{thm:finite-rep} represented $v$ as a sum of point masses.
The dual view
represents $v$ as a mixture of sharp
\emph{point-valuations}.  For each $p \in P$, let
$\delta_p$ be the valuation with $\delta_p(L) = 1$ if $p \in L$, and
$\delta_p(L) = 0$ otherwise:

\begin{theorem}[mixture characterization\leanref{valuation-eq-mix-deltapoint}]
\label{thm:mixture}
On a finite frame, every valuation $v$ satisfies
$v = \sum_p \mathrm{mass}(p)\cdot\delta_p$.
\end{theorem}

\noindent
Together with the fact that any mixture of valuations is again a
valuation\leanref{valuation-mix}, this theorem characterizes the
valuations on a finite frame exactly: they are the finite convex
mixtures of point-valuations, closed under mixing because modularity
is linear.  Section~\ref{sec:cox} shows modularity cannot instead be
\emph{derived} from more primitive axioms.  It must be posited
outright, and doing so is exactly what makes every valuation a
mixture of sharp, all-or-nothing certainties.

\subsection{Infinite Boundary Representation via Scott Continuity}

Both halves of the finite theorem fail for infinite frames:

\begin{example}[obstruction\leanref{exists-valuation-not-point-representable}]
\label{thm:obstruction}
On $\mathrm{Low}(\mathbb{N})$, the indicator $v$ of $\top$ is a
valuation with $\mathrm{mass}(p) = 0$ for every point $p$, and
$v(\top) = 1$.
\end{example}

\noindent
No point carries any of the mass.  The failure is one of
approximation: $v(\top)$ is not the supremum of $v(\pdown n)$
over the finite stages. In order to recover the representation,
we introduce Scott continuity:

\begin{definition}[directed set, Scott continuity]
A subset $D \subseteq \Omega$ is \emph{directed} if, for any finite
subset $F \subseteq D$, there exists an element $d \in D$ that is an
upper bound of $F$ (taking $F = \emptyset$ shows a directed set is
nonempty).  A
valuation $v$ is \emph{Scott-continuous} if
$v(\bigsqcup D) = \sup_{a \in D} v(a)$ for every directed $D$.
\end{definition}

The $\top$-indicator of Example~\ref{thm:obstruction} is not
Scott-continuous\leanref{topindicator-not-scott}, and Scott
continuity is exactly what restores the representation:

\begin{theorem}[recovery\leanref{eq-tsum-mass-of-scott}]
\label{thm:recovery}
Let $v$ be a valuation on $\mathrm{Low}(\mathbb{N})$. If
$v(\top) = \sup_n v(\pdown n)$
then $v(\top) = \sum_n \mathrm{mass}(n)$.
\end{theorem}

In general, on
any poset all of whose principal lower sets are finite, Scott
continuity implies zero diffuse part without additional preconditions.
However, subtraction in $\ENNReal$ is truncated at $0$, so if
$\sum_p \mathrm{mass}(p)$ were greater than $v(\top)$, the expression
$v(\top) - \sum_p \mathrm{mass}(p)$ would still evaluate to $0$.
Theorem~\ref{thm:decomposition} rules out
$\sum_p \mathrm{mass}(p) > v(\top)$ unconditionally, so the
subtraction below is never truncated:

\begin{theorem}[general decomposition\leanref{tsum-mass-le}]
\label{thm:decomposition}
For every valuation $v$ on $\mathrm{Low}(P)$ and an arbitrary poset $P$,
$\sum_p \mathrm{mass}(p) \le v(\top)$.
\end{theorem}

\noindent
By Theorem~\ref{thm:decomposition} the remainder
$v(\top) - \sum_p \mathrm{mass}(p)$ is a genuine difference.  Every valuation thus
splits unconditionally into an \emph{atomic} part carried by points
($\sum_p \mathrm{mass}(p)$, which is always a sub-probability), and a
\emph{diffuse} (nonnegative) remainder
$v(\top) - \sum_p \mathrm{mass}(p)$.
This is a constructive Lebesgue-style decomposition.  The finite and
Scott cases are the purely atomic extremes, and the $\top$-indicator
is the purely diffuse one.  The proof of
Theorem~\ref{thm:decomposition} re-runs the maximal-point induction on
each finite subset of $P$, bounding its partial sum by $v(\top)$,
then takes the supremum over all finite subsets to bound
$\sum_p \mathrm{mass}(p)$ itself.  
Its proof does not directly invoke an excluded-middle case split.
However, it relies on the (classical) decidability of the ordering
of the reals (a default in \textsf{mathlib}).
Representing the diffuse part on non-spatial frames remains the open
frontier.  The obstruction
there is again that the frame has \emph{too few points}, the same
one as behind the undecidability of Section~\ref{sec:bridges}.

\section{The Dempster--Shafer Reading}
\label{sec:ds}

Sections~\ref{sec:valuation}--\ref{sec:representation} independently
arrive at the same structure as Dempster--Shafer (DS)
theory~\cite{dempster1967,shafer1976}, including sub-additive complements, a
belief/plausibility interval, and mass functions.  
This section makes the identification exact in both statics and
dynamics.  In each case, it locates the precise point where excluded
middle makes a difference.

\subsection{Total Monotonicity via Constructive
Inclusion--Exclusion}

The $\neg\neg$-stable (\emph{regular}) elements of a frame form a
Boolean algebra with meet $\sqcap$ and join
$\neg\neg(a \sqcup b)$ (Glivenko).  Let us call the resulting
algebra the
\emph{Booleanization} of this set.

\begin{definition}[belief function, plausibility function]
\label{def:belief}
Let $B$ be a Boolean algebra.  A DS \emph{belief function} on $B$ is a
monotone map $\Bel : B \to [0,1]$ with $\Bel(\bot) = 0$ and
$\Bel(\top) = 1$ that is \emph{$n$-monotone} for
every $n \ge 2$: for all $A_1, \dots, A_n \in B$,
\[
  \Bel\Bigl(\bigvee_{i=1}^{n} A_i\Bigr) \;\ge\;
  \sum_{\emptyset \neq T \subseteq \{1,\dots,n\}} (-1)^{|T|+1}\,
  \Bel\Bigl(\bigwedge_{i \in T} A_i\Bigr).
\]
Its \emph{plausibility function} is the map $\Pl : B \to [0,1]$
given by $\Pl(A) = 1 - \Bel(\neg A)$.
\end{definition}

We show $v$, restricted
to its Booleanization, is such a belief function.  The
$2$-monotone case is immediate from modularity
\leanref{two-monotone}.  The full tower follows from the following
result:

\begin{theorem}[inclusion--exclusion\leanref{inclusion-exclusion}]
\label{thm:incl-excl}
For every valuation on every frame $\Omega$, every finite family
$a : \iota \to \Omega$, and every finite $s \subseteq \iota$,
inclusion--exclusion holds with
equality for the frame join:
\[
  v\Bigl(\,\bigsqcup_{i \in s} a_i\Bigr)
  \;=\; \sum_{\emptyset \neq T \subseteq s} (-1)^{|T|+1}\,
        v\Bigl(\,\bigsqcap_{i \in T} a_i\Bigr).
\]
\end{theorem}

We state this in Lean additively, with the odd-cardinality half of the
signed sum set against the even-cardinality half, so that no
$\ENNReal$ subtraction occurs.  The proof is by induction with
modularity and frame distributivity.

Replacing the frame join $\bigsqcup_i a_i$ on the equality's other
side with its double negation $\neg\neg\bigl(\bigsqcup_i a_i\bigr)$, the
Booleanization's own join, yields the following inequality:

\begin{theorem}[total monotonicity\leanref{infty-monotone}]
\label{thm:infty}
For every valuation on every frame $\Omega$, every finite family
$a : \iota \to \Omega$, and every finite $s \subseteq \iota$:
\[
  \sum_{\emptyset \neq T \subseteq s} (-1)^{|T|+1}
  v\Bigl(\bigsqcap_{i \in T} a_i\Bigr)
  \;\le\;
  v\Bigl(\neg\neg\bigl(\bigsqcup_{i \in s} a_i\bigr)\Bigr).
\]
\end{theorem}

The Booleanization's join is given by
$\neg\neg \circ \sqcup$, and its meets of regular
elements are again regular. Interpreted in this way, 
Theorem~\ref{thm:infty} states exactly
the $\infty$-monotone tower: $v$ restricted to its Booleanization is
a belief function.

The proof of Theorem~\ref{thm:infty} from
Theorem~\ref{thm:incl-excl} follows from the monotonicity along
$a \le \neg\neg a$.
Here, inclusion--exclusion does \emph{not} fail
constructively.  It is exact for the intuitionistic disjunction.
Applying the double-negation nucleus (a closure operator preserving
finite meets) $\neg\neg$ to the
join on the equality's right-hand side can only enlarge it. This 
follows from the fact that
$a \le \neg\neg a$ and $v$ are monotone, which turns the equality
into the inequality of Theorem~\ref{thm:infty}.  The amount by which
the right-hand side is enlarged is exactly the double-negation gap of
Section~\ref{sec:valuation}.

Definition~\ref{def:belief} is the classical notion, stated on an
abstract Boolean algebra.  In our setting each of its ingredients is
already present: the Boolean algebra is the Booleanization, the
belief function is $v$ itself, restricted to the regular elements and
evaluated against the Boolean join $\neg\neg \circ \sqcup$, and the
plausibility function is $\Pl(a) = 1 - v(\neg a)$.  Additive
probability does not become a belief function by weakening an axiom,
since modularity for $v$ is never relaxed.  What changes is only the
join used to state the tower.  The belief--plausibility interval is
governed by the slack:

\begin{theorem}[interval width\leanref{plausibility-sub-self}]
\label{thm:interval-width}
For every valuation $v$ on a frame $\Omega$ and every $a \in \Omega$,
$\Pl(a) - v(a) = \slack(a)$.
\end{theorem}


\subsection{The Conditioning Characterization}

DS theory carries two update rules that coincide classically.
\emph{Geometric} conditioning renormalizes belief, and under the DS
dictionary it \emph{is} the localic posterior
$v(a \sqcap b)/v(b)$ of Section~\ref{sec:valuation}.

\begin{definition}[Dempster's rule of conditioning\leanref{dempstercond}]
\label{def:dempster}
For a valuation $v$ on a frame $\Omega$, elements $a, b \in \Omega$, and
evidence of nonzero plausibility, $v(b) \neq 0$,
\emph{Dempster's rule of conditioning}~\cite{shafer1976}
states that
\[
  v(a \,\|\, b) \;=\;
  \frac{v(a \sqcup \neg b) - v(\neg b)}
       {1 - v(\neg b)}.
\]
The hypothesis $v(b) \neq 0$ makes the denominator nonzero, since
$v(b) + v(\neg b) \le 1$ gives $v(b) \le 1 - v(\neg b)$.
\end{definition}

\noindent
That is, the DS conditioning rule discards
the mass incompatible with the evidence and renormalizes by its
plausibility.
Both the geometric posterior of Section~\ref{sec:valuation} and
Dempster's rule of Definition~\ref{def:dempster} are again
valuations.  Their modularity
is the frame distributivity
$(a \sqcup \neg b) \sqcap (a' \sqcup \neg b)
= (a \sqcap a') \sqcup \neg b$.  Their disagreement is
given by the \emph{excluded-middle
gap} $\mathrm{emGap}(a,b) = v(a) - v(a \sqcap (b \sqcup
\neg b))$, i.e. the mass of $a$ outside $b$'s instance
of excluded middle:

\begin{lemma}[the join--complement identity\leanref{sup-compl-decomp}]
For every valuation $v$ on every frame $\Omega$ and every
$a, b \in \Omega$,
$v(a \sqcup \neg b) = v(\neg b) + v(a \sqcap b) +
\mathrm{emGap}(a,b)$.
\end{lemma}

Dempster's numerator therefore exceeds Bayes's by exactly
$\mathrm{emGap}(a,b)$.  This gap depends on $a$, but whether it
vanishes at \emph{every} $a$ depends only on
the slack at the evidence $b$:

\begin{theorem}[the conditioning characterization\leanref{dempstercond-eq-condval-iff-slack}]
\label{thm:cond-hinge}
For every valuation $v$ on every frame $\Omega$ and every
$b \in \Omega$ with $v(b) \neq 0$,
$\bigl(v(a \,\|\, b) = v(a \sqcap b)/v(b)$ for every
$a \in \Omega\bigr)$ iff $\slack(b) = 0$.
\end{theorem}

The hypothesis $\slack(b) = 0$ eliminates the gap
$\mathrm{emGap}(a,b)$ for every $a$ and makes the two normalizers
coincide, since $1 - v(\neg b) = v(b)$.
Conversely, testing at $a := b$, geometric conditioning gives
certainty $1$, while Dempster gives $v(b)/(1 - v(\neg b))$, which is
$1$ only if the slack vanishes.
Theorem~\ref{thm:cond-hinge} is the valuation-relative analogue of
Theorem~\ref{thm:hinge}: both are iff statements about the same
proposition, $b \sqcup \neg b = \top$.

We call the frame-level fact \emph{static} and the valuation-level
fact \emph{dynamic}: the former is an equation in the frame itself,
between $b \sqcup \neg b$ and $\top$, with no valuation involved,
while the latter tests the same proposition through $v$, since
$\slack(b) = 0$ is exactly $v(b \sqcup \neg b) = v(\top)$.  Dempster
and Bayes therefore coincide at $b$ whenever $v$ cannot distinguish
$b \sqcup \neg b$ from $\top$, even if the two differ in the frame
itself.  In the measure model, Dempster and Bayes disagree precisely
when the evidence's boundary carries mass.

\section{Cox's Theorem}
\label{sec:cox}

Cox's theorem~\cite{cox1946,jaynes2003} derives the probability
calculus from conditions on plausibility.  After Halpern's
counterexample~\cite{halpern1999}, its repaired
forms~\cite{vanhorn2003,arnborg-sjodin2000} fix the calculus up to
monotone regraduation.  Every known derivation is classical, and
Section~\ref{sec:hinge} showed why: the negation axiom R3 \emph{is}
excluded middle.  Dropping R3 leaves the product rule
unaffected, since it is logic-free, and requires the sum rule to become an axiom
rather than a theorem.  Separately, the naive constructive statement,
demanding a regraduation strictly monotone on all of $\mathbb{R}$, is
unsatisfiable (Section~\ref{sec:cox-theorem}).

\subsection{A Cox Structure without Negation}

\begin{definition}[Cox model\leanref{coxmodel}]
\label{def:coxmodel}
A \emph{Cox model} on a frame $\Omega$ is a plausibility assignment
$\mathrm{pl} : \Omega \times \Omega \to \mathbb{R}$, written
$\mathrm{pl}(a, b)$ for the plausibility of $a$ given $b$, together with
a \emph{conjunction functional} $F : \mathbb{R} \times \mathbb{R} \to
\mathbb{R}$, such that $\mathrm{pl}(\cdot, c)$ is monotone for every $c$,
the boundary conditions $\mathrm{pl}(\bot, c) = 0$ and
$\mathrm{pl}(\top, c) = 1$ hold for every $c$, the product rule
$\mathrm{pl}(a \sqcap b, c) = F(\mathrm{pl}(a, b \sqcap c),
\mathrm{pl}(b, c))$ holds, $F$ is associative, and $x \mapsto F(x, y)$ is
strictly monotone for every $y > 0$.
\end{definition}

The missing R3 axiom can be observed in the omission of
any constraint on the relation between
$\mathrm{pl}(\neg a, b)$ and $\mathrm{pl}(a, b)$.
The definition restricts the monotonicity of $F$ to
$0 < y$ to avoid a contradiction with the fact that
the boundary axioms force $F(x, 0) = 0$ for all
$x$.  An explicit model on the
non-classical chain $\ENNReal$ satisfies
this constraint\leanref{coxmodelennreal}, as $a = \top$ is provably decidable.

\subsection{The Sum Rule Is Irreducible}

Classically, additivity is \emph{derived}: R3 turns disjunction into
negated conjunction via De~Morgan, and the product rule does the
rest.  Constructively that route is closed, since
$a \vee b \neq \neg(\neg a \wedge \neg b)$ in general:

\begin{theorem}[the sum rule is irreducible\leanref{no-disjunction-functional}]
\label{thm:no-functional}
Let $\Omega$ be the five-element frame of lower sets of the poset
$o < a,\, o < b$ with $a, b$ incomparable.  Then there is a
monotone, normalized assignment $q$ on $\Omega$ that is additive on
disjoint joins yet not modular, and no binary operation $S$ on
values satisfies $q(x \sqcup y) = S(q(x), q(y))$ for all $x, y$.
Moreover, the
pairs $(\pdown a, \pdown a)$ and $(\pdown a,
\pdown b)$ have identical value-pairs $(\tfrac12, \tfrac12)$
but joins valued $\tfrac12$ and $1$.
\end{theorem}

\noindent
Unlike for a conjunction, a disjunction's plausibility cannot
generally be expressed as its disjuncts'
plausibilities via conditioning.  So, a constructive Cox
theorem \emph{must} posit the sum-rule half.  Modularity,
which says that the calculus is closed
under mixing certainties, is what
posits this, justified by Theorem~\ref{thm:mixture}.
The failure is not repaired by strengthening the logic short of
Boolean: the counterexample frame satisfies weak excluded
middle\leanref{isdemorgan-lowersetv}, since every nonempty lower set
contains the bottom point $o$, so the irreducibility persists under
that intermediate assumption.

\subsection{Cox Model Theorem}
\label{sec:cox-theorem}

A regraduation is a map $g : \mathbb{R} \to \ENNReal$ with
$g(0) = 0$.  No such map is strictly monotone on all of $\mathbb{R}$,
since strictness would force $g(-1) < g(0) = 0$, impossible in
$\ENNReal$.
The repair (strictness only on $[0,1]$, where plausibilities live)
is exactly the same domain restriction that Halpern's critique of
the classical theorem turns on.  The corrected statement is:

\begin{theorem}[the constructive Cox theorem\leanref{constructive-cox}]
\label{thm:cox}
For every frame $\Omega$, every Cox model on $\Omega$ whose
unconditional plausibility is modular
regraduates to a valuation: there exist $v : \mathrm{Valuation}\;
\Omega$ and $g$ strictly monotone on $[0,1]$ with $g(0)=0$,
$g(1)=1$, and for all
$a \in \Omega$, $v(a) = g(\mathrm{pl}(a, \top))$ .
\end{theorem}

The regraduation for the sum-rule half is straightforward
($g = \texttt{ENNReal.ofReal}$).  The product-rule half, showing $F$
regraduates to multiplication, rests on Acz\'el's associativity
theorem~\cite{aczel1966}, a statement of pure analysis whose type
never mentions $\Omega$.  \textsf{mathlib} has no Acz\'el theorem, no
H\"older-style embedding for ordered semigroups on intervals, and no
Cauchy functional equation under monotonicity, so we build the
analytic core from scratch, obtaining a normalized, strictly monotone
multiplicative generator on the positive
cone\leanref{exists-mul-generator}.  The construction also weakens
H\"older's classical hypotheses~\cite{hoelder1901}: the Archimedean
property is \emph{derived} from continuity and positivity rather than
assumed, and the interval form of divisibility follows from the
intermediate value theorem\leanref{exists-diag-eq} applied to a
plain divisibility hypothesis\leanref{scale}.  The
constructive/classical divide is therefore \emph{entirely} in the
sum rule: nothing in the product-rule apparatus mentions $\Omega$.
Appendix~\ref{app:aczel} gives the construction.  The residual
analytic gap is discussed in
Section~\ref{sec:conclusion}.

\section{Two Grounding Models}
\label{sec:bridges}

Two concrete models demonstrate that our abstract theory is 
non-degenerate.  One shows every
classical measure \emph{contains} an intuitionistic probability.
The other shows the intuitionistic theory cannot collapse back
without solving the halting problem.

\subsection{Measures through the Interior Operator}
\label{sec:interior-bridge}

G\"odel--McKinsey--Tarski interpret intuitionistic logic in modal
$\mathsf{S4}$: propositions are \emph{open} sets, and
the modal necessity operator $\Box$ of $\mathsf{S4}$ is interpreted
as the topological interior, written $\mathrm{int}(\cdot)$
below~\cite{godel1933,mckinsey-tarski1944}.
Our valuations quantify that same translation:

\begin{theorem}[measures as valuations\leanref{measure-tovaluationopens}]
\label{thm:measure-valuation}
Let $\mu$ be a probability measure on a space $X$, $A \subseteq X$,
and define $v(A) = \mu(\mathrm{int}\,A)$.  Since
$\mathrm{int}\,U = U$ for open $U$, $v$ restricts on the frame
$\mathcal{O}(X)$ to $U \mapsto \mu(U)$, and this restriction is a
valuation.
\end{theorem}

\noindent
Modularity is $\mu(U \cup V) + \mu(U \cap V) = \mu(U) + \mu(V)$.
The slack of an open $U$ is then $1 - \mu(U) -
\mu(\mathrm{int}\,U^{c})$, which is the measure of the
boundary $\partial U$.  We can now prove that, in general, slack is
exactly the boundary mass:

\begin{theorem}[slack is boundary mass\leanref{tovaluationopens-slack-eq-frontier}]
\label{thm:slack-frontier}
For any probability measure $\mu$ on a topological space $X$ and any
open $U$, $\slack(U) = \mu(\partial U)$.
\end{theorem}

\noindent
The Heyting complement of $U$ in $\mathcal{O}(X)$ is the interior of
its set complement. So, the undecided region, the complement of
$U \sqcup \neg U$, is exactly
$\overline{U} \setminus U = \partial U$.  The general sub-additivity
inequality of Section~\ref{sec:valuation} still holds here\leanref{interiormeasure-add-compl-le}.  This
identification of slack with boundary mass also exhibits the modal
identification $\Bel = P(\Box\,\cdot)$ of the DS
literature~\cite{ruspini1986} arising from the logic itself.
Theorem~\ref{thm:finite-rep} is the
converse in the finite case (in the general case, the converse is the open
problem, see Section~\ref{sec:representation}).

The same decomposition locates a failure of classical probability calculus
in conditioning, that is, in updating a belief about an open $A$ on
evidence given by an open $B$. The Boolean complement of $A$ splits into its
Heyting complement and its boundary,
$A^{c} = \mathrm{int}\,A^{c} \sqcup \partial A$, so a classical
posterior that treats $A$ as confirmed-or-refuted overstates the
refutation by the boundary mass the evidence retains:

\begin{theorem}[classical overconfidence in refutation\leanref{classical-posterior-compl-eq}]
\label{thm:classical-overconfidence}
Let $\mu$ be a probability measure on a topological space $X$, let
$v$ be its restriction to $\mathcal{O}(X)$, and let $A, B$ be opens
with $\mu(B) \neq 0$. Then
\[
  \frac{\mu(A \cap B)}{\mu(B)} = v(A \mid B)
  \qquad\text{and}\qquad
  \frac{\mu(A^{c} \cap B)}{\mu(B)}
  = v(\neg A \mid B) + \frac{\mu(\partial A \cap B)}{\mu(B)}.
\]
\end{theorem}

\noindent
The excess term is the conditional slack of $A$ given $B$, the mass
that the evidence leaves neither confirmed nor refuted. The classical
posterior does not discard this mass. It credits all of it to $A^{c}$,
an event no finite observation confirms. For the halting event of
Section~\ref{sec:computability-guard} with $B = \top$, the classical
credence in non-termination is $1 - p$ while the intuitionistic one is
$0$, and the difference $1 - p$ is exactly the slack.

\subsection{Universal Laws}
\label{sec:ravens}

Theorem~\ref{thm:classical-overconfidence} applies to universal laws such
as ``all ravens are black''~\cite{hempel1945}. Let $I$ be an infinite set
of objects and $O$ a topological space of kinds of object, with an open
set $\mathit{bad} \subseteq O$ of counterexamples, the non-black ravens. A
world is a function $x : I \to O$, and the worlds carry the product
topology, so an open set of worlds is a property verified by observing
finitely many objects. The law is
$H = \{x \mid x_i \notin \mathit{bad} \text{ for every } i\}$, and the
counterexample event is its complement
$C = \{x \mid x_i \in \mathit{bad} \text{ for some } i\}$, which is open.
For $j \in I$ and $A \subseteq O$ let $B_A = \{x \mid x_j \in A\}$, and
let $C_j$ be the event of a counterexample at an object other than $j$.

\begin{theorem}[universal laws\leanref{ravens}]
\label{thm:ravens}
Let $\mathit{bad}$ be nonempty.
(i) $\neg C = \bot$ in $\mathcal{O}(I \to O)$. So every valuation $v$
has $v(\neg C) = 0$, also after conditioning on any open, and
$\slack(C) = 1 - v(C)$.
(ii) For every probability measure $\mu$ on the worlds whose open sets
are measurable, $\mu(H) = \slack(C)$, and
$\mu(H \cap B)/\mu(B) = \mu(\partial C \cap B)/\mu(B)$ for every open
$B$ with $\mu(B) \ne 0$.
(iii) Let $A \subseteq O$ be open with $\mu(B_A) \ne 0$, and let $C_j$ be
independent of $B_A$ and of $B_{A \cap \mathit{bad}}$ under $\mu$. Then
\[
  \frac{\mu(C \cap B_A)}{\mu(B_A)}
  = \mu(C_j) + \frac{\mu(B_{A \cap \mathit{bad}})}{\mu(B_A)}\,
    \mu(C_j^{c}).
\]
\end{theorem}

\noindent
For (i), a neighbourhood of a world constrains finitely many objects, so
it contains a world with a counterexample at another object. Hence $C$
is dense and $H$ has empty interior. No finite observation verifies $H$,
and the intuitionistic credence in it is $0$ on any evidence. Part (ii)
follows from (i) and Theorems~\ref{thm:slack-frontier}
and~\ref{thm:classical-overconfidence}, since $\partial C = H$. The
classical credence in the law, prior or posterior, is boundary mass of
$C$, the mass of worlds the evidence has not refuted and cannot verify.
Part (iii) splits $C \cap B_A$ at object $j$. When $O$ is discrete,
$B_A$ is clopen, so $\slack(B_A) = 0$ and Dempster's rule gives the same
posteriors (Theorem~\ref{thm:cond-hinge}). The choice of update rule
plays no part below.

Hempel's paradox arises because $H$ is classically equivalent to ``every
non-black thing is a non-raven'', so a white sock found indoors is an
instance of $H$ and classically confirms it. In $\mathcal{O}(I \to O)$
the law is not an event: its intuitionistic counterpart $\neg C$ is
$\bot$, and evidence bears on it only through the open event $C$. Part
(iii) measures that bearing. Let $B_A$ say that object $j$ is a raven,
or that it is non-black. With $O$ discrete, observing its full kind, a black raven or a
white sock, is the case of $A$ a single kind outside $\mathit{bad}$,
and it gives the credence $\mu(C_j)$ in a counterexample in both cases.
The difference from the credence after selection alone is
$\mu(B_{A \cap \mathit{bad}})/\mu(B_A)$, the probability that the
selected object is a counterexample, times $\mu(C_j^{c})$. For a
selected raven the first factor is the probability that a raven is
non-black. For a selected non-black thing it is the probability that a
non-black thing is a raven, which is small when non-black things far
outnumber non-black ravens. So when $\mu(C_j^{c}) > 0$ the white sock
lowers the credence in a counterexample, by much less than the black
raven does. Both conditions matter. Without independence the direction
can reverse: Good gives a prior under which a black raven lowers the
credence in $H$~\cite{good1967}. And under an independent prior in which
every object is a counterexample with probability at least some
$\varepsilon > 0$, the second Borel--Cantelli lemma gives $\mu(C) = 1$,
so $\mu(H) = 0$ and no observation moves either. This is the
Bayesian resolution of Hosiasson-Lindenbaum and its successors, surveyed
by Fitelson and Hawthorne~\cite{hosiasson1940,fitelson-hawthorne2010}.
What the verification topology changes is the reading of the decrease
of $\mu(C)$. With finitely many objects and $O$ discrete, every set of
worlds is open, $\neg C = H$, and the two calculi coincide. With
infinitely many objects they differ. Classically, a decrease of $\mu(C)$
is an equal increase of the credence $\mu(H) = 1 - \mu(C)$ in the law.
Intuitionistically $v(\neg C)$ stays $0$ by (i), so the same decrease is
an increase of $\slack(C)$, the mass that is unrefuted and that no
evidence turns into verification. The resolution uses only the open
event $C$ and never treats $H$ as an event.

To characterize weak excluded middle spaces, we give the 
(standard) definitions required: 

\begin{definition}[extremally disconnected]
A topological space $X$ is \emph{extremally disconnected} if the
closure of every open set is again open.
\end{definition}

\begin{definition}[Stone space, Stonean space]
The \emph{Stone space} of a Boolean algebra $B$ is the space of
ultrafilters on $B$, topologized by the clopen basis $\{U : b \in U\}$
for $b \in B$. \emph{Stone duality}~\cite{johnstone1982} identifies
Boolean algebras with exactly the compact, Hausdorff, totally
disconnected spaces arising this way.  A \emph{Stonean space} is a compact Hausdorff space that
is also extremally disconnected. Equivalently, it is the Stone space
of a \emph{complete} Boolean algebra.
\end{definition}

\begin{remark}[weak excluded middle on opens\leanref{isdemorgan-opens-iff-extremallydisconnected}]
Section~\ref{sec:valuation} noted that imposing weak excluded middle
on the whole frame collapses slack to pure double-negation
instability.  For $\Omega = \mathcal{O}(X)$, this global condition
holds exactly when $X$ is extremally disconnected~\cite{johnstone1980}.
\end{remark}

An immediate corollary gives the characterization:

\begin{corollary}[Stonean characterization\leanref{isdemorgan-opens-iff-stonean}]
Given a compact Hausdorff space $X$, weak excluded middle on
$\mathcal{O}(X)$ holds exactly when $X$ is a Stonean space.
\end{corollary}

\subsection{Slack Does Not Determine the Valuation}
\label{sec:no-recovery}

Submodularity (Theorem~\ref{thm:slack-submodular}) is a necessary
condition on slack.  It is not sufficient: distinct valuations, each
satisfying every axiom a valuation must satisfy, can share a slack
function.

\begin{example}[boundaries in the three-point space\leanref{frontier-subset-ptz}]
\label{thm:threept-frontier}
Let $\mathrm{ThreePt} = \{p_A, p_B, p_Z\}$ carry the topology
generated by the two singletons $\{p_A\}$ and $\{p_B\}$, so its opens
are exactly $\emptyset$, $\{p_A\}$, $\{p_B\}$, $\{p_A, p_B\}$, and the
whole space.  For every open $U$, $\partial U \subseteq \{p_Z\}$.
\end{example}

\noindent
This is a fact about the topology alone, independent of any measure:
$p_Z$ has no proper open neighbourhood, so it is the only point any
open's frontier can contain.

\begin{example}[non-recovery\leanref{exists-ne-valuation-same-slack}]
\label{thm:no-recovery}
There exist valuations $v \ne v'$ on $\mathcal{O}(\mathrm{ThreePt})$ whose slack functions are identical: $\slack_v = \slack_{v'}$ as functions.
\end{example}

\noindent
For weights $a, b, c \ge 0$ with $a + b + c = 1$, the point measure
$\mu_{a,b,c} = a\,\delta_{p_A} + b\,\delta_{p_B} + c\,\delta_{p_Z}$,
read through Section~\ref{sec:interior-bridge}'s bridge, gives a
valuation $v_{a,b,c}$ on $\mathcal{O}(\mathrm{ThreePt})$.  By
Example~\ref{thm:threept-frontier} and
Theorem~\ref{thm:slack-frontier}, its slack at an open
$U$ is $\mu_{a,b,c}(\partial U)$, equal to $c$ when $p_Z \in \partial
U$ and to $0$ otherwise: a quantity depending only on $c$, not on how
the remaining mass splits between $a$ and $b$.  So $v_{1/4,1/4,1/2}$
and $v_{1/8,3/8,1/2}$ have identical slack functions while
disagreeing at $\{p_A\}$, where they take the values $1/4$ and $1/8$.
Both are fully modular, monotone, and normalized: nothing in
Theorem~\ref{thm:slack-submodular} or its proof distinguishes them.

\noindent
Weak excluded middle does not rescue recoverability.  If anything,
the fully classical case is the worst one for it.  When $X$ is
extremally disconnected, the De~Morgan gap vanishes
(Section~\ref{sec:interior-bridge}) and slack degenerates to the
double-negation gap alone, but on a space where excluded middle holds
outright that gap vanishes too, so slack is identically $0$ for
\emph{every} valuation.  The zero slack function is then
shared by every probability measure on $X$, which is the most extreme failure
of recovery available.

\subsection{The Computability Guard}
\label{sec:computability-guard}

The theory developed here is non-degenerate only if there are spaces
whose slack does not vanish identically.  Computability supplies
them:

\begin{definition}[Sierpi\'nski topology, semidecidable proposition]
\label{def:semidecidable}
The \emph{Sierpi\'nski space} $\mathbb{S}$ is the type \texttt{Prop} of
propositions carrying the topology with exactly three opens:
$\emptyset$, the set $\mathit{halts} = \{q \mid q\}$ of true
propositions, and $\mathbb{S}$
itself~\cite{escardo2004,rosolini1986}.  A proposition is
\emph{semidecidable} if it is an open in this topology.
\end{definition}

\noindent
The topology encodes verifiability.  A point of $\mathbb{S}$ is a
proposition, say ``this computation halts.''  The one nontrivial
observation available is membership in $\mathit{halts}$, established
by producing a proof (running the computation to termination).  No
open captures the complementary observation, and on the frame
$\mathcal{O}(\mathbb{S}) = \{\bot, \mathit{halts}, \top\}$ that
$\Sigma_1$ asymmetry is a theorem:

\begin{theorem}[the $\Sigma_1$ asymmetry\leanref{haltsopen-compl}]
\label{thm:sigma1-asymmetry}
$\neg\mathit{halts} = \bot$.
\end{theorem}

\noindent
Every nonempty open of $\mathbb{S}$ contains $\mathsf{True}$, so
$\bot$ is the only open disjoint from $\mathit{halts}$. Thus, no nonempty
open of $\mathbb{S}$ is evidence that the computation does not halt.

We now instantiate the measure-to-valuation construction of
Theorem~\ref{thm:measure-valuation} at $X = \mathbb{S}$,
allowing us to compute the slack:

\begin{example}[halting slack\leanref{slack-eq-boundary}]
\label{thm:slack-eq-boundary}
Let $\mu = p \cdot \delta_{\mathsf{True}} + (1-p) \cdot
\delta_{\mathsf{False}}$, and let $v$ be the resulting valuation.
Then
\[
  \slack(\mathit{halts}) \;=\; 1 - p \;=\;
  \mu\bigl(\partial\,\{q \mid q\}\bigr)
  \;=\; \mu\{\mathsf{False}\}.
\]
\end{example}

\noindent
Since $\mathit{halts}$ is already open, $v(\mathit{halts}) =
\mu(\mathit{halts}) = p$.  The DS ``ignorance'' of
Section~\ref{sec:ds} is here the mass of the single point
$\mathsf{False}$, the one outcome no finite observation can ever
resolve.  This probability is a fact about the model, not about
anyone's ignorance.

\paragraph{Chaitin's $\Omega$.}
The same bridge scales to algorithmic information theory.  Fix a
universal prefix-free machine $U$, \emph{prefix-free} in that no
halting program of $U$ is a proper prefix of another, and
\emph{universal} in that it simulates every other prefix-free
machine.  Give $2^{\omega}$, the space of infinite bit strings, the
fair-coin measure $\mu$.  The event $H$
of having a halting program of $U$ as a prefix is open, since
halting is established by a finite prefix
(Definition~\ref{def:semidecidable}), and prefix-freeness makes $H$
a disjoint union of cylinders, so
$v(H) = \mu(H) = \sum_{p \text{ halts}} 2^{-|p|}$: exactly Chaitin's
halting probability $\Omega$~\cite{chaitin1975}.  Its slack
$\mu(\partial H)$ is the mass of strings approximated arbitrarily
well by halting strings while never having a halting program as a
prefix, and $\Omega$'s classical uncomputability is the aggregate
form of the pointwise obstruction above.  The measure-theoretic core
is mechanized on Cantor space with \textsf{mathlib}'s infinite
product measure: basic cylinders are clopen with measure
$2^{-|w|}$\leanref{isopen-wordcylinder}, and a \emph{finite}
prefix-free set of words has pairwise disjoint cylinders of total
measure $\sum_w 2^{-|w|}$\leanref{cantormeasure-biunion-wordcylinder},
the Kraft-inequality content underlying $\Omega$.  The universal
machine itself, and the countable extension of the finite Kraft sum,
are computability theory beyond our scope.

\begin{definition}[halting problem]
For a code $c$ and input $n$, the \emph{halting problem} is the
predicate $H(c, n) = $ ``the machine with code $c$ halts on input
$n$,'' (\texttt{ComputablePred}). 
\end{definition}

\noindent
Slack cannot be eliminated by simply declaring a classical answer
for each code. The family of valuations
defined below attempts to do this:

\begin{definition}[sharp classical readout\leanref{sharpreadout-not-computable}]
For a given $n$, the \emph{sharp classical readout} is the family
of $\{0,1\}$-valued valuations $\rho_n(c) \in \mathrm{Valuation}(\mathcal{O}(\mathbb{S}))$,
with $\rho_n(c)(\mathit{halts}) = 1$ if $H(c,n)$
holds, and $\rho_n(c)(\mathit{halts}) = 0$ otherwise.
\end{definition}

\noindent
The readout is sharp but not slack-free. Since
$\neg\mathit{halts} = \bot$ (Theorem~\ref{thm:sigma1-asymmetry}), a
code that does not halt receives slack $1 - 0 = 1$, the maximum. Sharp
here means that the classical answer is known, not that the
complement rule holds on $\mathcal{O}(\mathbb{S})$.

\noindent
However, this family is provably not computable:

\begin{theorem}[the sharp readout is uncomputable\leanref{sharpreadout-not-computable}]
\label{thm:sharp-uncomputable}
For every $n$, the predicate $c \mapsto \bigl(\rho_n(c)(\mathit{halts}) = 1\bigr)$
is not computable.
\end{theorem}

\noindent
The proof identifies the predicate with $H(\cdot, n)$ and applies the
undecidability of the halting problem
(\texttt{ComputablePred.\allowbreak halting\_problem}).

\noindent
No route back to a classical, decidable treatment survives contact
with computability. Excluded middle guarantees that $H(c,n)$ is true
or false for any fixed $c$, but that guarantee is free: it names no
algorithm and exhibits no witness for which disjunct holds. Demanding
an actual witness, uniformly in $c$, is exactly $\rho_n$'s sharp
readout, and Theorem~\ref{thm:sharp-uncomputable} shows no algorithm
computes it. Aggregating over the whole family of halting programs,
rather than reading off one instance, does not help either:
$\Omega$ is a perfectly well-defined classical probability, yet
Chaitin showed it uncomputable by an independent argument turning on
algorithmic randomness~\cite{chaitin1975}. Neither the pointwise nor
the aggregate route gives classical probability a way to compute a
solution to the halting problem, and each is blocked by its own
uncomputability result.

\noindent
We can conclude that slack is not a defect of
the intuitionistic theory. On $\mathcal{O}(\mathbb{S})$, forcing the
complement rule at $\mathit{halts}$ leaves the single value
$v(\mathit{halts}) = 1$, so every code is assigned certain
termination. That assignment is computable but false of every
non-halting code. Any assignment that avoids this error, giving
$v(\mathit{halts}) < 1$ to each non-halting code, has positive slack
at every such code, and the one that is moreover sharp is the
uncomputable readout of Theorem~\ref{thm:sharp-uncomputable}.
For the sum rule, this guard plays the role the
explicit Cox-model instance plays for the product rule
(Section~\ref{sec:cox}): each axiom set in the paper comes with 
non-trivial models satisfying it.

\section{Product Valuations and Kernel Composition}
\label{sec:monad}

The structure
through which classical probability supports sequential statistical
reasoning is the Giry monad:

\begin{definition}[Giry monad~{\cite{giry1982}}]
\label{def:giry}
The \emph{Giry monad} $G$ sends a measurable space $X$ to the space
$G(X)$ of probability measures on $X$.  Its unit sends a point
$x \in X$ to the point mass at $x$, and its multiplication sends a
measure $M$ on $G(X)$ to the measure
$A \mapsto \int m(A) \,\mathrm{d}M(m)$.
\end{definition}

\noindent
Probabilistic reasoning then composes in stages, and, when the monad
is commutative, independent stages compose in either order.  The
question is whether the intuitionistic calculus supports the same
structure.

A \emph{frame
homomorphism} is a map preserving finite meets and arbitrary joins.
Pushforward along
a frame homomorphism $f : \Omega \to \Omega'$ sends a valuation $v$
on $\Omega'$ to $a \mapsto v(f(a))$ on $\Omega$, and this preserves
the valuation axioms, making $\mathrm{Val}(-)$ a functor on
locales, the localic counterpart of the functor part of the Giry
monad~\cite{giry1982}.  A frame homomorphism is the opposite of a 
locale map, so the functor laws are checked against
identity and composition of frame
homomorphisms\leanref{valuation-comap}.

\begin{definition}[valued frames\leanref{valf}]
\label{def:spval}
The \emph{valuation functor} $\mathrm{Val} : \mathrm{Frm}^{\mathrm{op}} \to
\mathrm{Set}$ sends a frame to its set of valuations and a frame
homomorphism $\varphi : P \to Q$ to $w \mapsto w \circ \varphi$. The
category $\mathrm{SpVal}$ of \emph{valued frames} is its category of
elements: objects are pairs $(P, v)$ with $v$ a valuation on $P$, and a
morphism $(P, v) \to (Q, w)$ is a frame homomorphism $\varphi : P \to Q$
with $w \circ \varphi = v$.
\end{definition}

\noindent
A morphism of valued frames is a measure-preserving map of locales in
the opposite direction.  The point valuations
$\delta_p$ (Section~\ref{sec:representation}) provide the candidate
(monad) unit.
Theorem~\ref{thm:finite-rep} shows that, on finite frames,
$\mathrm{Val}$ agrees (on objects) with the classical
finite-distribution monad on points.  
Flattening a valuation of valuations is
integration, and integration in two variables requires product
valuations.  A \emph{kernel} is a map assigning a valuation to each point
of a locale, giving a conditional distribution.  By Kock's
identification of commutative monads with Fubini
theorems~\cite{kock2012}, kernels compose independently of order
exactly when $\mathrm{Val}$ is a commutative monad.  

A valuation $u$ on a product frame is a \emph{product} of
valuations $v$ and $w$ when $u(A \times B) = v(A) \cdot w(B)$ for
all rectangles $A \times B$.  Let
$\mathrm{Low}(\mathbb{N} \times \mathbb{N})$ be the frame of lower
sets of the grid $\mathbb{N} \times \mathbb{N}$, ordered
component-wise.  On this frame, products are unique:

\begin{theorem}[product uniqueness\leanref{valuation-product-unique}]
\label{thm:product-unique}
Let $v, w$ be valuations on $\mathrm{Low}(\mathbb{N})$.  Any two
valuations on $\mathrm{Low}(\mathbb{N} \times \mathbb{N})$ that are
products of $v$ and $w$ are equal.
\end{theorem}

\noindent
By the inclusion--exclusion equality, two valuations agreeing on a
$\sqcap$-closed family agree on all finite joins of its
members\leanref{eq-on-finsetsup-of-eq-on}.  The rectangles are
closed under $\sqcap$, and in
$\mathrm{Low}(\mathbb{N} \times \mathbb{N})$ every lower set is a
\emph{finite} union of
rectangles\leanref{exists-rect-decomposition}, so agreement on
rectangles is agreement everywhere.  Valuations carry no continuity
axiom (Section~\ref{sec:representation}), 
so the $\sigma$-additivity driving the
classical uniqueness argument for product measures has no analogue
here.  The order structure of the frame does that work instead.

\begin{theorem}[product non-uniqueness\leanref{fubini-fails-for-valuations}]
\label{thm:product-not-unique}
On the powerset frame $\mathrm{Set}(\mathbb{N} \times \mathbb{N})$
there exist valuations $v, w$ and two \emph{distinct} products of
$v$ and $w$.
\end{theorem}

\noindent
An \emph{ultrafilter} on $\mathbb{N}$ is a collection of subsets,
closed under intersection and superset, containing exactly one of
$S$ and its complement for every $S \subseteq \mathbb{N}$.  It is
\emph{free} when it contains no finite set.  The map assigning $1$
to the members of an ultrafilter and $0$ to every other subset is a
valuation\leanref{ultrafiltervaluation}.  A rectangle $S \times T$
belongs to an ultrafilter on the product exactly when $S$ and $T$
belong to its images under the two coordinate
projections\leanref{rect-mem-iff}.  A free ultrafilter
$\mathcal{U}$ on $\mathbb{N}$ yields two ultrafilters on
$\mathbb{N} \times \mathbb{N}$.  A set belongs to the former when,
for $\mathcal{U}$-most first coordinates, its slice contains
$\mathcal{U}$-most second coordinates.  A set belongs to the latter
under the same test with the coordinates in the opposite order.

The two ultrafilters' valuations agree on
\emph{every} rectangle yet differ on the upper triangle
$\{(m, m') \mid m < m'\}$\leanref{product-valuation-not-unique}.
Integrating the same marginal data in the two orders thus
disagrees.  Fubini fails for valuations, as it does for non-additive
integration~\cite{denneberg1994,ghirardato1997}.  By Kock's
correspondence, no canonical commutative
multiplication exists for unrestricted valuations.  This half of the
dichotomy is not specific to intuitionism, since powerset-frame
valuations are exactly the classical finitely additive probability
measures~\cite{bhaskara-rao1983}. Dropping continuity is what admits
the non-uniqueness. For continuous valuations on continuous domains the
product is unique and Fubini's theorem holds~\cite{jones1990}, and they
form the probabilistic powerdomain~\cite{jones-plotkin1989}.  A monad therefore requires restricting either
the frames or the valuations.

A first step in the second direction is to restrict the mixing
measures to countable ones. The components of a mixture are arbitrary
valuations, so this restricts how valuations are combined, not which
valuations occur. With countable weights the order of mixing does not
matter, by the interchange of double series. In
Theorem~\ref{thm:product-not-unique} the mixing measures are ultrafilter
valuations, which are not countable mixtures:

\begin{proposition}[countable kernel composition\leanref{mixcountable}]
\label{prop:mixcountable}
\begin{enumerate}
\item A countable mixture of valuations is a
  valuation\leanref{mixcountable}.
\item A mixture of mixtures flattens to a mixture with product
  weights, which is discrete
  Chapman--Kolmogorov\leanref{mixcountable-mixcountable}.
\item The one-point mixture is the
  identity\leanref{mixcountable-unique}.
\item Iterated countable mixing is
  order-independent\leanref{mixcountable-comm}.
\end{enumerate}
\end{proposition}

\noindent
Restricted to countable-mixture presentations, the data of a monad are
all in place: the category of locales, $\mathrm{Val}$ restricted to
countable mixtures as the functor, the point valuations as the unit,
and countable mixing as the multiplication.
Proposition~\ref{prop:mixcountable} verifies exactly the laws this
needs: item~1 keeps the multiplication inside the fragment, item~2 is
its associativity, item~3 its unit law, and item~4 its commutativity.
What is not yet settled is whether this descends to valuations
themselves, since a valuation can admit more than one countable-mixture
presentation, and the diffuse part of Section~\ref{sec:representation}
is not captured by any countable presentation at all.

\section{Valuations Beyond Frames}
\label{sec:lattice}

Definition~\ref{def:frame} fixes the event algebra to be a frame: a
complete lattice in which finite meets distribute over arbitrary
joins $\bigsqcup_{i \in I} a_i$, the operation Heyting implication
$\Rightarrow$ (the right adjoint to $a \sqcap (-)$) is built from.
Definition~\ref{def:valuation}'s four axioms, normalization at
$\bot$ and $\top$, monotonicity, and modularity
$v(a) + v(b) = v(a \sqcup b) + v(a \sqcap b)$, use only the
\emph{binary} join $\sqcup$; none of the four mentions the
\emph{arbitrary} join $\bigsqcup_{i \in I}$. Completeness is used
only where later sections evaluate a valuation on an infinite join,
as in the Scott-continuity results of
Section~\ref{sec:representation}, or build $\Rightarrow$ and $\neg$
from it, as in Section~\ref{sec:hinge}'s slack.
Section~\ref{sec:conclusion} leaves open how much of the
construction is tied to frames specifically, and whether some
ambient category sends an event algebra to a probability calculus
in general. This section settles both questions for one fragment of
the construction, the sharp/point correspondence of
Section~\ref{sec:hinge}: it isolates the exact algebraic hypothesis
that correspondence needs and shows the hypothesis cannot be
dropped. It does not recover $\Rightarrow$ or $\neg$ without
completeness: the adjoint defining $\Rightarrow$ is built from an
arbitrary join of witnesses, so it is unavailable once completeness
goes, even if only finite joins are missing. The valuations studied
below are negation-free by construction, not frame valuations with
completeness quietly removed.

\paragraph{The functor without completeness.}
A \emph{bounded lattice} is a lattice with a least element $\bot$
and a greatest element $\top$; a \emph{bounded lattice homomorphism}
preserves $\bot$, $\top$, $\sqcup$, and $\sqcap$.
Definition~\ref{def:valuation}'s axioms make sense over any bounded
lattice $\Omega$, complete or not:

\begin{definition}[lattice valuation]
\label{def:lvaluation}
A \emph{lattice valuation} on a bounded lattice $\Omega$ is a
monotone function $v : \Omega \to \ENNReal$ with $v(\bot) = 0$,
$v(\top) = 1$, and, for all $a, b \in \Omega$,
$v(a) + v(b) = v(a \sqcup b) + v(a \sqcap b)$.
\end{definition}

\begin{lstlisting}
structure LValuation (Ω : Type*) [Lattice Ω] [BoundedOrder Ω] where
  toFun : Ω → ℝ≥0∞
  map_bot' : toFun ⊥ = 0
  map_top' : toFun ⊤ = 1
  mono'    : Monotone toFun
  modular' : ∀ a b, toFun a + toFun b
             = toFun (a ⊔ b) + toFun (a ⊓ b)
\end{lstlisting}

\noindent
A lattice valuation on a frame is exactly a valuation in the sense
of Definition~\ref{def:valuation}; the two notions differ only in
how much structure the ambient $\Omega$ carries.

Section~\ref{sec:monad} pulls a valuation back along a frame
homomorphism and checks the assignment is functorial (identity,
composition). That argument pattern-matches on $\sqcup$, $\sqcap$,
$\bot$, $\top$ alone and never touches $\bigsqcup$:

\begin{proposition}[pullback without completeness]
\label{prop:lvaluation-comap}
For bounded lattices $\Gamma, \Omega$, a bounded lattice
homomorphism $f : \Gamma \to \Omega$, and a lattice valuation $v$ on
$\Omega$, $v \circ f$ is a lattice valuation on $\Gamma$. This
assignment sends the identity homomorphism to the identity and
composition of homomorphisms to composition of the induced maps.
\end{proposition}

\noindent
(Mechanized as \texttt{LValuation.comap}, \texttt{comap\_id}, and
\texttt{comap\_comp} in \texttt{GeneralValuation.\allowbreak lean}.)
$\mathrm{Val}(-)$ is therefore a functor on bounded lattices with
bounded lattice homomorphisms, of which the frame-level functor of
Section~\ref{sec:monad} is the restriction to the subcategory of
frames and frame homomorphisms. Completeness contributes nothing to
functoriality.

\paragraph{Separation needs only distributivity.}
Theorem~\ref{thm:sharp-iff-point} characterizes what a sharp
valuation \emph{is}, a complement-indicator of a prime ideal, but
not whether there are enough of them to distinguish elements of
$\Omega$. That is a separate fact:

\begin{theorem}[sharp valuations separate points]
\label{thm:sharp-separating}
For every frame $\Omega$ and $a, b \in \Omega$ with $a \not\le b$,
there is a sharp valuation $v$ with $v(a) = 1$ and $v(b) = 0$.
\end{theorem}

\noindent
(Mechanized as \texttt{exists\_sharp\_separating} in
\texttt{Points.lean}.) The proof runs the same prime-ideal
separation theorem that built the counterexample of
Theorem~\ref{thm:hinge}
(\texttt{DistribLattice.prime\_ideal\_of\_disjoint\_filter\_ideal}),
on the principal filter of $a$ and the principal ideal of $b$
rather than on $a \sqcup \neg a$ and $\top$. As a corollary, any two
distinct $a, b \in \Omega$ are distinguished by some sharp
valuation: $\Omega$ need not be linear, so $a \ne b$ alone does not
hand back $a \le b$ or $b \le a$, but it forces at least one failure
of $\le$, and Theorem~\ref{thm:sharp-separating} applies in whichever
direction fails. This is the cash-out of the first qualifying
condition of Section~\ref{sec:conclusion}: the certain valuations do
not merely characterize a single point of the locale in isolation,
together they reconstruct the order on $\Omega$. Reconstructing the
order is not the same as reconstructing the full Heyting structure:
$\neg$ and $\Rightarrow$, as just noted, are adjoints built from an
arbitrary join and carry no meaning over a plain distributive
lattice. What this theorem recovers is the order-theoretic core of
the first qualifying condition, not the negation fragment of
intuitionistic logic, which stays frame-specific.

Tracing the proof of Theorem~\ref{thm:sharp-separating} shows it
never invokes $\bigsqcup$; the only place frame structure enters is
\texttt{DistribLattice.prime\_ideal\_of\_disjoint\_filter\_ideal}
itself, which needs distributivity and nothing beyond it:

\begin{theorem}[separation over distributive lattices]
\label{thm:lvaluation-separating}
For every distributive bounded lattice $\Omega$ and $a, b \in
\Omega$ with $a \not\le b$, there is a sharp lattice valuation $v$
on $\Omega$ with $v(a) = 1$ and $v(b) = 0$.
\end{theorem}

\noindent
(Mechanized as \texttt{LValuation.exists\_sharp\_separating} in
\texttt{GeneralValuation.lean}.) Theorem~\ref{thm:lvaluation-separating}
reproves Theorem~\ref{thm:sharp-separating} with the hypothesis
\texttt{Order.Frame} weakened to \texttt{DistribLattice}, and
nothing else about the argument changes. The condition doing the
work in the sharp/point correspondence is distributivity;
completeness is not load-bearing for either half of it.

\paragraph{Distributivity is not a proof artifact.}
The natural question is whether
Theorem~\ref{thm:lvaluation-separating}'s hypothesis is a genuine
dividing line or an artifact of routing the argument through a
mathlib lemma named for distributive lattices. It is genuine, and
the failure it guards against is total, not merely a failure of
separation:

\begin{definition}[the diamond]
\label{def:m3}
$M_3$ is the five-element lattice with a bottom $z$, a top $o$, and
three pairwise-incomparable atoms $a, b, c$, with $x \sqcup y = o$
and $x \sqcap y = z$ for every two distinct atoms $x, y$.
\end{definition}

\noindent
$M_3$ is the smallest lattice that is not distributive:
$a \sqcap (b \sqcup c) = a \sqcap o = a$, while
$(a \sqcap b) \sqcup (a \sqcap c) = z \sqcup z = z$. It is modular,
so this failure is specific to distributivity rather than some
coarser defect of the order.

\begin{theorem}[no sharp valuation on the diamond]
\label{thm:m3-no-sharp}
No lattice valuation on $M_3$ is sharp.
\end{theorem}

\noindent
(Mechanized as \texttt{M3.no\_sharp\_\allowbreak valuation} in
\texttt{NonDistributive.\allowbreak lean}, together with
\texttt{M3.not\_\allowbreak distributive} for the displayed failure of
distributivity above.) Modularity applied to each of the three
pairs of distinct atoms forces $v(x) + v(y) = 1$, since
$x \sqcup y = \top$ and $x \sqcap y = \bot$. If $v$ were sharp,
$v(a), v(b), v(c) \in \{0, 1\}$ would have to satisfy all three
pairwise sum-to-one equations at once, and a two-element codomain
cannot support that: fixing any one value propagates through the
other two equations to a contradiction with the third. The collapse
on $M_3$ is therefore not a failure of separation, points that
happen to coincide, but a failure of existence: there is no
all-or-nothing valuation on $M_3$ to separate anything with, and the
first qualifying condition of Section~\ref{sec:conclusion} is
unsatisfiable outright.

\paragraph{Where this leaves the boundary.}
Frames are complete distributive lattices.
Theorem~\ref{thm:lvaluation-separating} shows completeness is not
needed for the sharp/point correspondence's separation half, and
Theorem~\ref{thm:m3-no-sharp} shows distributivity cannot be given
up in exchange. The event algebras on which sharp valuations reconstruct the order,
the order-theoretic core of the first qualifying condition of
Section~\ref{sec:conclusion}, are exactly the distributive bounded
lattices, no more and no less. The second qualifying condition and
the negation-dependent machinery of
Sections~\ref{sec:hinge}--\ref{sec:bridges}, slack, the complement
bound, the De Morgan gap, are untouched by this section and remain
tied to frames: nothing here shows, or claims, that $\neg$ or
$\Rightarrow$ survives dropping completeness. This bears on the
orthomodular event lattices of quantum probability discussed in
Section~\ref{sec:related}: giving up distributivity, as those
lattices do, does not merely change which valuations exist, it can
remove sharp valuations altogether, so the construction of this
paper does not, without further modification, have anything to say
about that setting even at the order-theoretic level examined here.
Whether some weaker notion of certainty survives on non-distributive
lattices, in place of the $\{0,1\}$-valued sharp valuations examined
here, is not addressed by this section.

\section{Related Work}
\label{sec:related}

\paragraph{Informal precedents.}
Pairing probability with intuitionistic logic is not a new idea,
though the literature on it is scattered.  Weatherson develops
intuitionistic probability axiomatically and defends it with
Dutch-book arguments~\cite{weatherson2003}.  His axioms match our
valuation axioms exactly, with
negation derived rather than assumed, so the object itself and the
bound $v(a) + v(\neg a) \le 1$ are his.  What is new here
is the decomposition of that bound's slack into $\mathrm{dnGap}$ and
$\mathrm{deMorganGap}$ and its determinacy from $v$ alone
(Theorem~\ref{thm:slack-explicit}).  
Paris shows that the assignments avoiding a Dutch book are the mixtures
of a logic's two-valued valuations, and that when the valuations respect
$\wedge$ and $\vee$, as in intuitionistic and modal logics, they are
exactly the normalized, monotone, modular
assignments~\cite{paris2001}. On finite frames this is
Theorem~\ref{thm:mixture}. Williams obtains the same set from accuracy
domination, under the name generalized probabilism~\cite{williams2012}.
Paris also shows, from Shafer's results, that the mixtures of the
``known true'' valuations, which make a sentence true when it follows
from what is known, are exactly the Dempster--Shafer belief
functions~\cite{paris1994,paris2001,shafer1976}. These valuations do not
respect $\vee$.  We mechanize an analogous fact for our
intuitionistic setting: the $2$-monotone case follows directly from
modularity\leanref{two-monotone}, and the full tower is
Theorem~\ref{thm:infty}.  The
identification $\Bel = P(\Box\,\cdot)$ connecting belief
functions~\cite{dempster1967,shafer1976} to modal logic is
established~\cite{ruspini1986}.  Section~\ref{sec:bridges} shows this
identification is not a separate axiom for us but a special case of
our frame-level slack, recovered automatically once the frame is a
topology and $\neg$ is read as the interior of the complement
\leanref{tovaluationopens-slack-eq-frontier}.  

On the localic
side, valuations on
frames are standard objects of constructive measure theory, with
Vickers's work and the Coquand--Spitters development as canonical
references~\cite{vickers2008,coquand-spitters2009}.  These normally
build in Scott continuity and target additive, Riesz-style goals.
Section~\ref{sec:representation} shows Scott continuity to be a
substantive assumption (it decides point-representability on the
frames we examine), so our base structure omits it.  Cox, Jaynes,
Halpern's counterexample, and the Van Horn and Arnborg--Sj\"odin
repairs~\cite{cox1946,jaynes2003,halpern1999,vanhorn2003,
arnborg-sjodin2000} are classical throughout, and Cox's and Van
Horn's derivations both lean on double-negation elimination.
Section~\ref{sec:cox} proves
the sum-rule half without that move, at the price of positing
modularity as an axiom rather than deriving it
(Example~\ref{thm:no-functional}).  Van Horn observes
that dropping axiom R3 yields an interval-valued theory, exactly the
belief/plausibility pair our slack separates.  To our knowledge, no
prior work, informal or formal, states either characterization
theorem: complement rule $\iff$ excluded middle
(Theorem~\ref{thm:hinge}), or Dempster $=$ Bayes $\iff$ zero slack
at the evidence (Theorem~\ref{thm:cond-hinge}).

\paragraph{Prior formalizations.}
Two bodies of mechanized work are adjacent, and neither overlaps.
ALEA~\cite{audebaud-paulin2009} formalizes a discrete distribution
monad in Coq for randomized programs: additive probability, no
event logic in our sense.  Closest is the
HoTT/Coq line of Faissole and Spitters on \emph{valuations} over
lower reals~\cite{faissole-spitters2017}, done constructively with
Scott continuity and aimed at integration (Riesz) and a Giry-style
valuation \emph{monad}.  Their valuations differ from ours in
carrying Scott continuity, and their monad is the Scott-continuous
instance of the question examined in
Section~\ref{sec:monad}.  Our development drops the continuity
their base structure builds in, and the two mechanizations do not
otherwise overlap.  On the domain-theoretic side, commutativity of
the full valuations monad on the category of directed-complete
partial orders is open.  A hierarchy of commutative submonads,
generated from simple valuations (finite mixtures of point
valuations), has been constructed
instead~\cite{jia2021commutative}.  
Our countably-presented fragment, the valuations given as countable
mixtures (Proposition~\ref{prop:mixcountable}), is the frame-side
analogue.

\paragraph{Non-Boolean probability elsewhere.}
Quantum probability on orthomodular lattices, with Gleason's
theorem~\cite{birkhoff-vonneumann1936,gleason1957}, is the
long-standing precedent for re-basing probability on a non-Boolean
event logic.  Orthomodular lattices give up distributivity itself,
while our frames stay distributive by definition
(Definition~\ref{def:frame}), locating all the non-classical
behavior in the failure of excluded middle alone.
Section~\ref{sec:quantum} applies the construction from certain states
to a quantum event logic.  



Linear logic
marks a different kind of departure.  Its event algebras are
quantales, complete lattices equipped with an associative tensor
$\otimes$ that need not be idempotent~\cite{girard1987}.  The
lattice reduct of a quantale is already an instance of
Section~\ref{sec:lattice}'s generalization, but $\otimes$ itself is
not $\sqcap$, and a valuation grading it cannot be phrased in
Definition~\ref{def:lvaluation}'s vocabulary without first deciding
what modularity even means for a non-idempotent operation.
Section~\ref{sec:recipe} replaces modularity by the affine laws valid on
the certain states.  The probabilistic semantics that does exist for
linear logic, Girard's probabilistic coherence spaces as developed by
Danos and Ehrhard~\cite{danos-ehrhard2011}, grades the multiplicative
structure directly rather than valuing a lattice.
Section~\ref{sec:recipe} shows that it contains the calculus of
certain states at every multiplicative--additive formula and equals it
at implications between data types, and Section~\ref{sec:second} that
the inclusion is strict at second order.  



Free probability,
fuzzy measures, and imprecise
probability~\cite{walley1991} vary other parameters (independence,
additivity, single-valuedness).  Where Walley's coherent lower
previsions posit the belief/plausibility gap directly as a
coherence primitive, our slack (Section~\ref{sec:valuation}) is
derived, forced by modularity and negation rather than assumed.
Fritz's synthetic Markov-category
framework~\cite{fritz2020markov} varies the ambient structure
categorically rather than order-theoretically, axiomatizing
conditional independence and sufficient statistics abstractly
across discrete, measure-theoretic, and Gaussian probability alike.
It remains unmechanized and is orthogonal to our variation in the
underlying logic.  Wolpert and Kinney model
mathematics itself as a stochastic process, replacing binary
theoremhood with a probability distribution~\cite{wolpert-kinney2024}.
Theoremhood is a semidecidable event, so the computability guard of
Section~\ref{sec:bridges} bears directly on that program:
sharp credences about such events are uncomputable.  Our
variation is specifically the
\emph{logic}, holding the Cox-style methodology fixed.  Within DS
theory, the two conditionings and their inequivalence are
classical~\cite{shafer1976,fagin-halpern1991}.  What is new in
Theorem~\ref{thm:cond-hinge} is the exact logical characterization
of their agreement, and its mechanization.

\section{Conclusion}
\label{sec:conclusion}

Many properties of programs are only semidecidable, so the Boolean
event algebra Kolmogorov's axioms~\cite{kolmogorov1933} presuppose is
not a good fit for reasoning about them.  The notion of graded 
belief does not presuppose Boolean logic.
Classicality is imposed only via the complement rule
(Theorem~\ref{thm:hinge}).  Its dynamic counterpart is the agreement
of Dempster and Bayesian conditioning (Theorem~\ref{thm:cond-hinge}).  
Everything else
survives the intuitionistic substrate intact: conditioning, Bayes
symmetry, disjoint additivity, total probability over genuine
partitions, and inclusion--exclusion itself
(Theorem~\ref{thm:incl-excl}).  The weaker logic adds the slack
structure, along with its two-gap decomposition, and the
atomic/diffuse divide. For this new structure,
we demonstrate a belief-function reading through
$\neg\neg$, and give a computability-theoretic reason that structure
cannot be dispensed with: forcing it away on a semidecidable event
is exactly as hard as solving the halting problem
(Section~\ref{sec:bridges}).

\paragraph{Functoriality and the monad question.}
For a fixed choice of logic, the
connection to the categorical picture is concrete.  Pushforward along
a frame homomorphism $f : \Omega \to \Omega'$ sends a valuation $v$
on $\Omega'$ to $(f_* v)(a) = v(f(a))$ on $\Omega$, and this
preserves all four valuation axioms, making $\mathrm{Val}(-)$ a
functor on locales, the localic counterpart of the functor part of
the Giry monad~\cite{giry1982}.  This is mechanized end to end, the
pushforward together with the two functor laws checked against
identity and composition of frame homomorphisms\leanref{valuation-comap}.  The
point valuations $\delta_p$ of Section~\ref{sec:representation} are
the unit, already formalized, and Theorem~\ref{thm:finite-rep} shows
that on finite frames $\mathrm{Val}$ agrees, on objects, with the
classical finite-distribution monad on points.  The multiplication
is where the substance is.  Flattening a valuation of valuations is
integration, hence product valuations and Fubini, and by Kock's
identification of commutative monads with Fubini
theorems~\cite{kock2012}, whether statistics-as-kernel-composition
transports to a weaker logic is precisely the question whether that
logic's valuation functor is a commutative monad.  Two
machine-checked results delimit the multiplication from both sides.

On the chain product, the question of \emph{which} product valuation
to take does not arise, since the frame
$\mathrm{Low}(\mathbb{N} \times \mathbb{N})$ is \emph{rigid}: every
lower set is a \emph{finite} union of rectangles, because beyond a
stabilization height every occupied column is full, so the full
columns and finitely many rows account for everything
\leanref{exists-rect-decomposition}.  The rigidity mechanism is
general.  By the inclusion--exclusion equality, two valuations
agreeing on any $\sqcap$-closed family agree on all finite joins of
its members\leanref{eq-on-finsetsup-of-eq-on}, and the chain
product is the instance where that family is the rectangles.  The
product valuation there, whenever it exists, is therefore unique,
with no continuity assumed\leanref{valuation-product-unique}.

On the powerset frame $\mathrm{Set}(\mathbb{N} \times \mathbb{N})$,
whose elements are far from finite unions of rectangles,
non-uniqueness is a theorem.  The indicator of an ultrafilter is a
valuation\leanref{ultrafiltervaluation}, and any ultrafilter on a
product splits rectangles through its marginal pushforwards
\leanref{rect-mem-iff}.  Taking a free ultrafilter $\mathcal{U}$
on $\mathbb{N}$, the two iteration orders of its tensor with itself
and its diagonal pushforward carry identical values on \emph{every}
rectangle, yet they differ on the upper triangle
$\{(m, m') \mid m < m'\}$
\leanref{product-valuation-not-unique}.  The two
orders of iterated integration of the same marginal data therefore
disagree: Fubini fails, concretely
\leanref{fubini-fails-for-valuations}.  Compare the Fubini
pathologies of non-additive
integration~\cite{denneberg1994,ghirardato1997}.  By Kock's
correspondence, no canonical commutative multiplication exists for
unrestricted valuations.  On the powerset frame these valuations are
finitely additive probability charges in the sense
of~\cite{bhaskara-rao1983}, so this half of the dichotomy is not
specific to intuitionism: the freedom is already present in the
classical finitely additive setting, and dropping continuity is what
admits it.  A monad therefore requires restricting either the
frames or the valuations.

\paragraph{The obstruction to a monad.}
The question upstream of the multiplication is worth stating
precisely, though what follows is an assessment rather than a
theorem.  A monad needs a category for $\mathrm{Val}$ to be an
endofunctor on.  The space of valuations must itself be an object, a
``space'' whose opens one can speak of.  For Scott-continuous
valuations that space is a locale and monad structure is available,
constructively, through the probabilistic powerdomain
line~\cite{jones-plotkin1989,vickers2008,coquand-spitters2009}.
Even there, exhibiting a cartesian closed category closed under the
construction is the Jung--Tix problem, open since 1998 for the
classical probabilistic powerdomain functor
itself~\cite{jung-tix1998}, though reframed versions of the category
have recently been settled.  Dropping continuity makes the
difficulty prior rather than merely worse, since the standard
presentations of the valuation space use continuity essentially.

One candidate route to a well-posed question does exist.  The
localic treatments present the space of continuous valuations as the
classifying locale of a propositional geometric theory~\cite{vickers2008,
coquand-spitters2009}, with generators asserting $q < v(a)$ for
rational $q$ and axioms for monotonicity, normalization, modularity
(expressible with rational parameters and countable disjunctions),
and continuity.  Deleting the continuity axiom leaves a geometric
theory, so a classifying locale still exists on general grounds, and
classically its points are exactly the valuations of this paper.  We
have verified none of the structure this candidate would need.  The
unit looks unproblematic, but the multiplication requires
integrating the evaluation maps against a merely modular,
discontinuous valuation, where the sum of two lower-real-valued maps
is an infinite join, and the diffuse part of
Example~\ref{thm:obstruction} is precisely what such joins fail to
capture.  We offer the formulation as a
target, not a result.

A more modest strategy, with a better track record, is to carve out
a fragment rather than construct the whole monad.  On the
measure-theoretic side, the failure of cartesian closure for
measurable spaces was answered by quasi-Borel
spaces~\cite{heunen2017qbs}.  On the domain-theoretic side,
commutativity of the full valuations monad on \textsf{DCPO} is open,
and the response is a hierarchy of commutative \emph{submonads}
generated from simple valuations~\cite{jia2021commutative}.  Our
decomposition theorem names the corresponding fragment here: the
atomic part.  As a first machine-checked step in that direction, we
show that the countably-presented fragment composes correctly with
no space structure needed at all.  A countable mixture of valuations
is a valuation\leanref{mixcountable}, a mixture of mixtures
flattens to a mixture with product weights, which is discrete
Chapman--Kolmogorov\leanref{mixcountable-mixcountable}, the
one-point mixture is the identity\leanref{mixcountable-unique},
and iterated countable mixing is order-independent\leanref{mixcountable-comm}, all by the unconditional Fubini of
$\ENNReal$ sums over sigma types.  This is associative, unital,
commutative kernel composition for countably-atomic data, in direct
contrast with the non-uniqueness above.  It is not yet a monad: the
presentations are not canonical, we do not quotient by ``presents the
same valuation,'' and the diffuse part is exactly what no countable
presentation captures.

In summary, the unit, the finite level, and countable kernel
composition with its commutativity are proved.  Rigidity on chain
products and non-uniqueness on the powerset frame are proved, so the
general multiplication is settled: it is not canonically determined.
What remains is the ambient category, a formulation problem before
it is anything else.

\paragraph{Open problems and future work.}
Four problems are open: (1)~representing the diffuse part on
non-spatial frames, (2)~completing the Acz\'el development
(continuity of the generator and its extension beyond the cone),
(3)~deriving modularity from a criterion more primitive than the
mixture characterization (Theorem~\ref{thm:mixture}), knowing it
cannot come from disjunction data alone
(Theorem~\ref{thm:no-functional}), and (4)~a machine-checked
Dutch-book theorem in the style of~\cite{weatherson2003}, which to
our knowledge no proof assistant has for \emph{any} logic.  On the
monad, two continuations are concrete: quotient the countable
presentations into a Kleisli category of countably-atomic kernels,
and mechanize the classifying-locale formulation above, which needs
classifying locales of geometric theories, infrastructure
\textsf{mathlib} currently lacks.  Both are confined to a fragment,
since the unrestricted multiplication is provably not canonical.
Beyond intuitionistic logic there is no ambient category in which
``logic $\mapsto$ its probability theory'' is a functor.  What
transfers is the method: state the classical desiderata, delete the
axioms the target logic no longer licenses, and let a proof
assistant check what survives, what splits, and what becomes
unsatisfiable.  Linear logic is the natural next case: its event
algebras are quantales, a valuation there should grade the
consumption of evidence, and the analogues of our representation
theorems and of slack, for the failure of contraction, are
open~\cite{danos-ehrhard2011}.


\begin{acks}
Claude (Anthropic) was used to assist with drafting and debugging
the Lean formalization.
\end{acks}

\bibliographystyle{ACM-Reference-Format}
\bibliography{refs}

\appendix
\section{Values in the Extended Reals}
\label{app:ennreal}

Valuations take values in the extended non-negative reals $\ENNReal$,
following the measure-theoretic convention of \textsf{mathlib}.  
As a result,
sums indexed by an arbitrary family are unconditionally defined, with no
summability hypothesis, and they commute over dependent sums: for a
function $f$ on a sigma type $\Sigma\,a,\,\beta(a)$, one has
$\sum_{p} f(p) = \sum_{a} \sum_{b} f(a, b)$ with no side condition.  This
is what lets the countable-mixture results of
Section~\ref{sec:monad} hold as stated.  Additionally, subtraction is
truncated at $0$.  This never affects the main text, since every
subtraction there is applied to a pair already ordered by monotonicity,
as with the mass of Section~\ref{sec:representation}, so the truncation
is never triggered.

\section{Lean Statements}
\label{app:lean}

\subsection{Introduction}

\subsection{Valuations on Frames}

\phantomsection\label{lean:valuation}
\begin{lstlisting}
structure Valuation (Ω : Type*) [Order.Frame Ω] where
  toFun : Ω → ℝ≥0∞
  map_bot' : toFun ⊥ = 0
  map_top' : toFun ⊤ = 1
  mono' : Monotone toFun
  modular' : ∀ a b, toFun a + toFun b = toFun (a ⊔ b) + toFun (a ⊓ b)
\end{lstlisting}

\phantomsection\label{lean:add-compl-le-one}
\begin{lstlisting}
theorem Valuation.add_compl_le_one (v : Valuation Ω) (a : Ω) :
    v a + v aᶜ ≤ 1
\end{lstlisting}

\phantomsection\label{lean:valuation-add-compl-eq-sup}
\begin{lstlisting}
theorem Valuation.add_compl_eq_sup (v : Valuation Ω) (a : Ω) :
    v a + v aᶜ = v (a ⊔ aᶜ)
\end{lstlisting}

\phantomsection\label{lean:valuation-le-one}
\begin{lstlisting}
theorem Valuation.le_one (v : Valuation Ω) (a : Ω) : v a ≤ 1
\end{lstlisting}

\phantomsection\label{lean:valuation-slack-eq-dngap-add-demorgangap}
\begin{lstlisting}
theorem Valuation.slack_eq_dnGap_add_deMorganGap (v : Valuation Ω) (a : Ω) :
    v.slack a = v.dnGap a + v.deMorganGap a
\end{lstlisting}

\phantomsection\label{lean:valuation-slack-eq-sub-add-sub}
\begin{lstlisting}
theorem Valuation.slack_eq_sub_add_sub (v : Valuation Ω) (a : Ω) :
    v.slack a = (v aᶜᶜ - v a) + (1 - (v aᶜᶜ + v aᶜ))

theorem Valuation.le_compl_compl (v : Valuation Ω) (a : Ω) : v a ≤ v aᶜᶜ
\end{lstlisting}

\phantomsection\label{lean:valuation-slack-submodular}
\begin{lstlisting}
theorem Valuation.slack_submodular (v : Valuation Ω) (a b : Ω) :
    v.slack (a ⊔ b) + v.slack (a ⊓ b) ≤ v.slack a + v.slack b
\end{lstlisting}

\phantomsection\label{lean:isdemorgan}
\begin{lstlisting}
def Order.Frame.IsDeMorgan (Ω : Type*) [Order.Frame Ω] : Prop :=
  ∀ a : Ω, aᶜᶜ ⊔ aᶜ = ⊤
\end{lstlisting}

\phantomsection\label{lean:valuation-slack-eq-dngap-of-isdemorgan}
\begin{lstlisting}
theorem Valuation.slack_eq_dnGap_of_isDeMorgan
    (hΩ : Order.Frame.IsDeMorgan Ω) (v : Valuation Ω) (a : Ω) :
    v.slack a = v.dnGap a
\end{lstlisting}

\phantomsection\label{lean:condval-mul}
\begin{lstlisting}
noncomputable def Valuation.condVal (v : Valuation Ω) (b : Ω) (hb : v b ≠ 0) :
    Valuation Ω

theorem Valuation.condVal_mul (v : Valuation Ω) (b : Ω) (hb : v b ≠ 0) (a : Ω) :
    v.condVal b hb a * v b = v (a ⊓ b)

theorem Valuation.condVal_symm (v : Valuation Ω) {a b : Ω}
    (ha : v a ≠ 0) (hb : v b ≠ 0) :
    v.condVal b hb a * v b = v.condVal a ha b * v a
\end{lstlisting}

\phantomsection\label{lean:cond-add-compl-le}
\begin{lstlisting}
theorem Valuation.cond_add_compl_le (v : Valuation Ω) (a b : Ω) :
    v (a ⊓ b) + v (aᶜ ⊓ b) ≤ v b
\end{lstlisting}

\phantomsection\label{lean:add-compl-eq-one-of-complemented}
\begin{lstlisting}
theorem Valuation.add_compl_eq_one_of_complemented (v : Valuation Ω) {a : Ω}
    (ha : a ⊔ aᶜ = ⊤) : v a + v aᶜ = 1
\end{lstlisting}

\phantomsection\label{lean:classical-additivity}
\begin{lstlisting}
variable [CompleteBooleanAlgebra Ω]

theorem classical_additivity (v : Valuation Ω) (a : Ω) :
    v a + v aᶜ = 1
\end{lstlisting}

\phantomsection\label{lean:classical-slack-zero}
\begin{lstlisting}
variable [CompleteBooleanAlgebra Ω]

theorem classical_slack_zero (v : Valuation Ω) (a : Ω) :
    v.slack a = 0
\end{lstlisting}

\phantomsection\label{lean:valuation-mix}
\begin{lstlisting}
noncomputable def Valuation.mix (w : ι → ℝ≥0∞) (hw : ∑ i, w i = 1)
    (v : ι → Valuation Ω) : Valuation Ω
\end{lstlisting}

\subsection{Where Classicality Lives}

\phantomsection\label{lean:hinge}
\begin{lstlisting}
def Valuation.HasClassicalNegation (v : Valuation Ω) : Prop :=
  ∀ a, v a + v aᶜ = 1

theorem hasClassicalNegation_of_em (hem : ∀ a : Ω, a ⊔ aᶜ = ⊤)
    (v : Valuation Ω) : v.HasClassicalNegation

theorem em_of_forall_hasClassicalNegation
    (h : ∀ v : Valuation Ω, v.HasClassicalNegation) : ∀ a : Ω, a ⊔ aᶜ = ⊤
\end{lstlisting}

\phantomsection\label{lean:sharp-iff-point}
\begin{lstlisting}
theorem sharp_iff_point (v : Valuation Ω) :
    v.IsSharp ↔ ∃ (J : Order.Ideal Ω) (hJ : J.IsPrime) (htop : ⊤ ∉ J),
      v = Ideal.toValuation J hJ htop
\end{lstlisting}

\phantomsection\label{lean:sharp-scott-iff-completelyprimepoint}
\begin{lstlisting}
theorem sharp_scott_iff_completelyPrimePoint (v : Valuation Ω) :
    v.IsSharp ∧ v.ScottContinuous ↔
      ∃ (J : Order.Ideal Ω) (hJ : Ideal.IsCompletelyPrime J) (htop : ⊤ ∉ J),
        v = Ideal.toValuation J hJ.1 htop
\end{lstlisting}

\subsection{Valuation Representation}

\phantomsection\label{lean:valuation-mass}
\begin{lstlisting}
noncomputable def Valuation.mass (v : Valuation (LowerSet P)) (p : P) : ℝ≥0∞
\end{lstlisting}

\phantomsection\label{lean:valuation-eq-sum-mass}
\begin{lstlisting}
theorem Valuation.eq_sum_mass (v : Valuation (LowerSet P)) (U : LowerSet P) :
    v U = ∑ p ∈ U.toFinset, v.mass p

theorem Valuation.sum_mass (v : Valuation (LowerSet P)) : ∑ p, v.mass p = 1
\end{lstlisting}

\phantomsection\label{lean:valuation-topmf}
\begin{lstlisting}
noncomputable def Valuation.toPMF (v : Valuation (LowerSet P)) : PMF P
\end{lstlisting}

\phantomsection\label{lean:valuation-exists-pmf-presentation}
\begin{lstlisting}
theorem Valuation.exists_pmf_presentation {Θ : Type u} [Order.Frame Θ] [Finite Θ]
    (v : Valuation Θ) :
    ∃ (P : Type u) (_ : PartialOrder P) (_ : Fintype P) (_ : DecidableEq P)
      (e : Θ ≃o LowerSet P),
      (∀ a, v a = ∑ p ∈ (e a).toFinset, (v.mapEquiv e).mass p) ∧
        ∑ p, (v.mapEquiv e).mass p = 1
\end{lstlisting}

\phantomsection\label{lean:valuation-eq-mix-deltapoint}
\begin{lstlisting}
noncomputable def deltaPoint (p : P) : Valuation (LowerSet P)

theorem Valuation.eq_mix_deltaPoint (v : Valuation (LowerSet P)) :
    v = Valuation.mix v.mass v.sum_mass (fun p => deltaPoint p)
\end{lstlisting}

\phantomsection\label{lean:exists-valuation-not-point-representable}
\begin{lstlisting}
theorem exists_valuation_not_point_representable :
    ∃ v : Valuation (LowerSet ℕ), ∑' p, v.mass p ≠ v ⊤

theorem topIndicator_mass (n : ℕ) : topIndicator.mass n = 0
\end{lstlisting}

\phantomsection\label{lean:topindicator-not-scott}
\begin{lstlisting}
theorem topIndicator_not_scott :
    (⨆ n, topIndicator (LowerSet.Iic n)) ≠ topIndicator ⊤
\end{lstlisting}

\phantomsection\label{lean:eq-tsum-mass-of-scott}
\begin{lstlisting}
theorem eq_tsum_mass_of_scott (v : Valuation (LowerSet ℕ))
    (hsc : v ⊤ = ⨆ n, v (LowerSet.Iic n)) : v ⊤ = ∑' n, v.mass n

variable [∀ p : P, Finite (Set.Iic p)]

theorem Valuation.isPurelyAtomic_of_scott (v : Valuation (LowerSet P))
    (hsc : v ⊤ = ⨆ (S : Finset P), v (S.sup LowerSet.Iic)) : IsPurelyAtomic v
\end{lstlisting}

\phantomsection\label{lean:tsum-mass-le}
\begin{lstlisting}
theorem tsum_mass_le (v : Valuation (LowerSet P)) : ∑' p, v.mass p ≤ v ⊤
\end{lstlisting}

\subsection{The Dempster--Shafer Reading}

\phantomsection\label{lean:inclusion-exclusion}
\begin{lstlisting}
theorem Valuation.inclusion_exclusion (v : Valuation Ω) (s : Finset ι)
    (a : ι → Ω) : v (s.sup a) + v.iesEven s a = v.iesOdd s a
\end{lstlisting}

\phantomsection\label{lean:infty-monotone}
\begin{lstlisting}
theorem Valuation.infty_monotone (v : Valuation Ω) (s : Finset ι) (a : ι → Ω) :
    v.iesOdd s a ≤ v ((s.sup a)ᶜᶜ) + v.iesEven s a
\end{lstlisting}

\phantomsection\label{lean:two-monotone}
\begin{lstlisting}
theorem Valuation.two_monotone (v : Valuation Ω) (a b : Ω) :
    v a + v b ≤ v (a ⊔ b)ᶜᶜ + v (a ⊓ b)
\end{lstlisting}

\phantomsection\label{lean:plausibility-sub-self}
\begin{lstlisting}
noncomputable def Valuation.plausibility (v : Valuation Ω) (a : Ω) : ℝ≥0∞ := 1 - v aᶜ

theorem Valuation.plausibility_sub_self (v : Valuation Ω) (a : Ω) :
    v.plausibility a - v a = v.slack a
\end{lstlisting}

\phantomsection\label{lean:sup-compl-decomp}
\begin{lstlisting}
theorem Valuation.sup_compl_decomp (v : Valuation Ω) (a b : Ω) :
    v (a ⊔ bᶜ) = v bᶜ + (v (a ⊓ b) + v.emGap a b)
\end{lstlisting}

\phantomsection\label{lean:dempstercond}
\begin{lstlisting}
noncomputable def Valuation.dempsterCond (v : Valuation Ω) (b : Ω)
    (hb : v b ≠ 0) : Valuation Ω

theorem Valuation.dempsterCond_apply (v : Valuation Ω) (b : Ω) (hb : v b ≠ 0)
    (a : Ω) : v.dempsterCond b hb a = (v (a ⊔ bᶜ) - v bᶜ) / (1 - v bᶜ)
\end{lstlisting}

\phantomsection\label{lean:dempstercond-eq-condval-iff-slack}
\begin{lstlisting}
theorem Valuation.dempsterCond_eq_condVal_iff_slack (v : Valuation Ω) {b : Ω}
    (hb : v b ≠ 0) :
    (∀ a, v.dempsterCond b hb a = v.condVal b hb a) ↔ v.slack b = 0
\end{lstlisting}

\subsection{Cox's Theorem}

\phantomsection\label{lean:coxmodel}
\begin{lstlisting}
structure CoxModel (Ω : Type*) [Order.Frame Ω] where
  pl : Ω → Ω → ℝ
  mono_left : ∀ c, Monotone fun a => pl a c
  pl_bot : ∀ c, pl ⊥ c = 0
  pl_top : ∀ c, pl ⊤ c = 1
  F : ℝ → ℝ → ℝ
  conj : ∀ a b c, pl (a ⊓ b) c = F (pl a (b ⊓ c)) (pl b c)
  F_assoc : ∀ x y z, F (F x y) z = F x (F y z)
  F_strictMono_left : ∀ y, 0 < y → StrictMono fun x => F x y
\end{lstlisting}

\phantomsection\label{lean:isdemorgan-lowersetv}
\begin{lstlisting}
theorem isDeMorgan_lowerSetV : Order.Frame.IsDeMorgan (LowerSet V)
\end{lstlisting}

\phantomsection\label{lean:no-disjunction-functional}
\begin{lstlisting}
theorem modularity_irreducible :
    ∃ (Ω : Type) (_ : Order.Frame Ω) (q : Ω → ℝ≥0∞),
      Monotone q ∧ q ⊥ = 0 ∧ q ⊤ = 1 ∧
      (∀ x y : Ω, x ⊓ y = ⊥ → q (x ⊔ y) = q x + q y) ∧
      ¬ (∀ x y : Ω, q x + q y = q (x ⊔ y) + q (x ⊓ y))

theorem no_disjunction_functional :
    ¬ ∃ S : ℝ≥0∞ → ℝ≥0∞ → ℝ≥0∞,
        ∀ x y : LowerSet V, qV (x ⊔ y) = S (qV x) (qV y)
\end{lstlisting}

\phantomsection\label{lean:constructive-cox}
\begin{lstlisting}
theorem constructive_cox (M : CoxModel Ω)
    (hmod : ∀ a b : Ω, M.pl a ⊤ + M.pl b ⊤ = M.pl (a ⊔ b) ⊤ + M.pl (a ⊓ b) ⊤) :
    ∃ (v : Valuation Ω) (g : ℝ → ℝ≥0∞),
      StrictMonoOn g (Set.Icc 0 1) ∧ g 0 = 0 ∧ g 1 = 1 ∧
        ∀ a, v a = g (M.pl a ⊤)
\end{lstlisting}

\phantomsection\label{lean:coxmodelennreal}
\begin{lstlisting}
noncomputable def coxModelENNReal : CoxModel ℝ≥0∞
\end{lstlisting}

\phantomsection\label{lean:exists-mul-generator}
\begin{lstlisting}
theorem Scale.exists_mul_generator :
    ∃ G : ℝ → ℝ,
      (∀ x y, S.u ≤ x → S.u ≤ y → G (S.F x y) = G x * G y) ∧
      (∀ x y, S.u ≤ x → x < y → G x < G y)

theorem Scale.aczelStatement_cone :
    ∃ G : ℝ → ℝ,
      (∀ x y, S.u ≤ x → S.u ≤ y → G (S.F x y) = G x * G y) ∧
      StrictMonoOn G (Set.Ici S.u)
\end{lstlisting}

\phantomsection\label{lean:exists-diag-eq}
\begin{lstlisting}
theorem exists_diag_eq {F : ℝ → ℝ → ℝ} {a b c : ℝ} (hab : a ≤ b)
    (hcont : ContinuousOn (fun x => F x x) (Set.Icc a b))
    (hc : c ∈ Set.Icc (F a a) (F b b)) : ∃ x ∈ Set.Icc a b, F x x = c
\end{lstlisting}

\phantomsection\label{lean:scale}
\begin{lstlisting}
structure Scale where
  F : ℝ → ℝ → ℝ
  hdiv : ∀ c, ∃ x, F x x = c
  hassoc : ∀ x y z, F (F x y) z = F x (F y z)
  hpos : ∀ c x, x < F x c
  hmono : ∀ x₁ x₂ y₁ y₂ : ℝ, x₁ ≤ x₂ → y₁ ≤ y₂ → F x₁ y₁ ≤ F x₂ y₂
  hcont : Continuous fun p : ℝ × ℝ => F p.1 p.2
  hbot : ∀ x, Filter.Tendsto (fun z => F x z) Filter.atBot (nhds x)
  u : ℝ
\end{lstlisting}

\subsection{Two Grounding Models}

\phantomsection\label{lean:measure-tovaluationopens}
\begin{lstlisting}
noncomputable def MeasureTheory.Measure.toValuationOpens
    (μ : Measure X) [IsProbabilityMeasure μ] : Valuation (Opens X)
\end{lstlisting}

\phantomsection\label{lean:tovaluationopens-slack-eq-frontier}
\begin{lstlisting}
theorem toValuationOpens_slack_eq_frontier (μ : Measure X)
    [IsProbabilityMeasure μ] (U : Opens X) :
    μ.toValuationOpens.slack U = μ (frontier (U : Set X))
\end{lstlisting}

\phantomsection\label{lean:classical-posterior-compl-eq}
\begin{lstlisting}
theorem classical_posterior_compl_eq (μ : Measure X)
    [IsProbabilityMeasure μ] (a b : Opens X)
    (hb : μ.toValuationOpens b ≠ 0) :
    μ ((a : Set X)ᶜ ∩ b) / μ b
      = μ.toValuationOpens.condVal b hb aᶜ
        + μ (frontier (a : Set X) ∩ b) / μ b
\end{lstlisting}

\phantomsection\label{lean:interiormeasure-add-compl-le}
\begin{lstlisting}
theorem interiorMeasure_add_compl_le (μ : Measure X) [IsProbabilityMeasure μ]
    (S : Set X) : μ.interiorMeasure S + μ.interiorMeasure Sᶜ ≤ 1
\end{lstlisting}

\phantomsection\label{lean:isdemorgan-opens-iff-extremallydisconnected}
\begin{lstlisting}
theorem isDeMorgan_opens_iff_extremallyDisconnected :
    Order.Frame.IsDeMorgan (Opens X) ↔ ExtremallyDisconnected X
\end{lstlisting}

\phantomsection\label{lean:isdemorgan-opens-iff-stonean}
\begin{lstlisting}
theorem isDeMorgan_opens_iff_stonean [CompactSpace X] [T2Space X] :
    Order.Frame.IsDeMorgan (Opens X) ↔ ExtremallyDisconnected X
\end{lstlisting}

\phantomsection\label{lean:frontier-subset-ptz}
\begin{lstlisting}
inductive ThreePt : Type
  | ptA : ThreePt
  | ptB : ThreePt
  | ptZ : ThreePt

theorem frontier_subset_ptZ {U : Set ThreePt} (hU : IsOpen U) :
    frontier U ⊆ {ptZ}
\end{lstlisting}

\phantomsection\label{lean:exists-ne-valuation-same-slack}
\begin{lstlisting}
theorem exists_ne_valuation_same_slack :
    ∃ v v' : Valuation (Opens ThreePt), v.slack = v'.slack ∧ v ≠ v'
\end{lstlisting}

\phantomsection\label{lean:haltsopen-compl}
\begin{lstlisting}
theorem haltsOpen_compl : haltsOpenᶜ = (⊥ : Opens Prop)
\end{lstlisting}

\phantomsection\label{lean:slack-eq-boundary}
\begin{lstlisting}
theorem slack_eq_boundary {p : ℝ≥0∞} (hp1 : p ≤ 1) :
    (sierpinskiValuation hp1).slack haltsOpen
      = coinMeasure p (frontier {q : Prop | q})
\end{lstlisting}

\phantomsection\label{lean:sharpreadout-not-computable}
\begin{lstlisting}
noncomputable def haltingWeight (n : ℕ) (c : Code) : ℝ≥0∞ :=
  if (c.eval n).Dom then 1 else 0

noncomputable def sharpReadout (n : ℕ) (c : Code) : Valuation (Opens Prop) :=
  sierpinskiValuation (haltingWeight_le_one n c)

theorem sharpReadout_not_computable (n : ℕ) :
    ¬ComputablePred fun c : Code => sharpReadout n c haltsOpen = 1
\end{lstlisting}

\phantomsection\label{lean:isopen-wordcylinder}
\begin{lstlisting}
theorem isOpen_wordCylinder (w : List Bool) : IsOpen (wordCylinder w)

theorem cantorMeasure_wordCylinder (w : List Bool) :
    cantorMeasure (wordCylinder w) = (1 / 2) ^ w.length
\end{lstlisting}

\phantomsection\label{lean:cantormeasure-biunion-wordcylinder}
\begin{lstlisting}
theorem disjoint_wordCylinder_of_not_prefix {v w : List Bool}
    (hv : ¬ v <+: w) (hw : ¬ w <+: v) :
    Disjoint (wordCylinder v) (wordCylinder w)

theorem cantorMeasure_biUnion_wordCylinder {F : Finset (List Bool)}
    (hF : PrefixFree F) :
    cantorMeasure (⋃ w ∈ F, wordCylinder w) = ∑ w ∈ F, (1 / 2 : ℝ≥0∞) ^ w.length
\end{lstlisting}

\subsection{The Monad Question}

\phantomsection\label{lean:valuation-comap}
\begin{lstlisting}
def Valuation.comap (f : FrameHom Γ Ω) (v : Valuation Ω) : Valuation Γ

theorem Valuation.comap_id (v : Valuation Ω) : v.comap (FrameHom.id Ω) = v

theorem Valuation.comap_comp (f : FrameHom Γ Ω) (g : FrameHom Δ Γ)
    (v : Valuation Ω) : v.comap (f.comp g) = (v.comap f).comap g
\end{lstlisting}

\phantomsection\label{lean:exists-rect-decomposition}
\begin{lstlisting}
theorem exists_rect_decomposition (W : LowerSet (ℕ × ℕ)) :
    ∃ (N : ℕ) (A B : ℕ → LowerSet ℕ),
      W = (Finset.range N).sup fun i => A i ×ˢ B i
\end{lstlisting}

\phantomsection\label{lean:eq-on-finsetsup-of-eq-on}
\begin{lstlisting}
theorem Valuation.eq_on_finsetSup_of_eq_on {G : Set Ω}
    (hG : ∀ g ∈ G, ∀ g' ∈ G, g ⊓ g' ∈ G) (u u' : Valuation Ω)
    (h : ∀ g ∈ G, u g = u' g) {ι : Type*} [DecidableEq ι]
    (s : Finset ι) (a : ι → Ω) (ha : ∀ i ∈ s, a i ∈ G) :
    u (s.sup a) = u' (s.sup a)
\end{lstlisting}

\phantomsection\label{lean:valuation-product-unique}
\begin{lstlisting}
theorem Valuation.product_unique (v w : Valuation (LowerSet ℕ))
    {u u' : Valuation (LowerSet (ℕ × ℕ))}
    (hu : ∀ A B : LowerSet ℕ, u (A ×ˢ B) = v A * w B)
    (hu' : ∀ A B : LowerSet ℕ, u' (A ×ˢ B) = v A * w B) : u = u'
\end{lstlisting}

\phantomsection\label{lean:ultrafiltervaluation}
\begin{lstlisting}
noncomputable def ultrafilterValuation (W : Ultrafilter ι) : Valuation (Set ι)
\end{lstlisting}

\phantomsection\label{lean:rect-mem-iff}
\begin{lstlisting}
theorem rect_mem_iff {α β : Type*} {W : Ultrafilter (α × β)}
    {S : Set α} {T : Set β} :
    S ×ˢ T ∈ W ↔ S ∈ W.map Prod.fst ∧ T ∈ W.map Prod.snd
\end{lstlisting}

\phantomsection\label{lean:product-valuation-not-unique}
\begin{lstlisting}
theorem product_valuation_not_unique :
    ∃ u u' : Valuation (Set (ℕ × ℕ)),
      (∀ S T : Set ℕ, u (S ×ˢ T) = u' (S ×ˢ T)) ∧ u ≠ u'
\end{lstlisting}

\phantomsection\label{lean:fubini-fails-for-valuations}
\begin{lstlisting}
theorem fubini_fails_for_valuations :
    ∃ (v w : Valuation (Set ℕ)) (u u' : Valuation (Set (ℕ × ℕ))),
      (∀ S T : Set ℕ, u (S ×ˢ T) = v S * w T) ∧
      (∀ S T : Set ℕ, u' (S ×ˢ T) = v S * w T) ∧ u ≠ u'
\end{lstlisting}

\phantomsection\label{lean:mixcountable}
\begin{lstlisting}
noncomputable def Valuation.mixCountable {ι : Type*} (w : ι → ℝ≥0∞)
    (hw : ∑' i, w i = 1) (v : ι → Valuation Ω) : Valuation Ω
\end{lstlisting}

\phantomsection\label{lean:mixcountable-mixcountable}
\begin{lstlisting}
theorem Valuation.mixCountable_mixCountable {ι : Type*} {κ : ι → Type*}
    (w : ι → ℝ≥0∞) (hw : ∑' i, w i = 1)
    (u : ∀ i, κ i → ℝ≥0∞) (hu : ∀ i, ∑' j, u i j = 1)
    (v : ∀ i, κ i → Valuation Ω) :
    Valuation.mixCountable w hw
        (fun i => Valuation.mixCountable (u i) (hu i) (v i))
      = Valuation.mixCountable (fun p : Σ i, κ i => w p.1 * u p.1 p.2)
          (tsum_sigma_mul_eq_one w hw u hu) (fun p => v p.1 p.2)
\end{lstlisting}

\phantomsection\label{lean:mixcountable-unique}
\begin{lstlisting}
theorem Valuation.mixCountable_unique (v : Valuation Ω) :
    Valuation.mixCountable (fun _ : PUnit => 1) (by simp) (fun _ => v) = v
\end{lstlisting}

\phantomsection\label{lean:mixcountable-comm}
\begin{lstlisting}
theorem Valuation.mixCountable_comm {ι κ : Type*}
    (w : ι → ℝ≥0∞) (hw : ∑' i, w i = 1) (u : κ → ℝ≥0∞) (hu : ∑' j, u j = 1)
    (v : ι → κ → Valuation Ω) :
    Valuation.mixCountable w hw (fun i => Valuation.mixCountable u hu (v i))
      = Valuation.mixCountable u hu
          (fun j => Valuation.mixCountable w hw (fun i => v i j))
\end{lstlisting}

\subsection{Valuations Beyond Frames}

\noindent Theorem~\ref{thm:sharp-separating}.
\phantomsection\label{lean:sharp-separating}
\begin{lstlisting}
theorem exists_sharp_separating {a b : Ω} (hab : ¬ a ≤ b) :
    ∃ v : Valuation Ω, v.IsSharp ∧ v a = 1 ∧ v b = 0
\end{lstlisting}

\noindent Definition~\ref{def:lvaluation}.
\phantomsection\label{lean:lvaluation}
\begin{lstlisting}
structure LValuation (Ω : Type*) [Lattice Ω] [BoundedOrder Ω] where
  toFun : Ω → ℝ≥0∞
  map_bot' : toFun ⊥ = 0
  map_top' : toFun ⊤ = 1
  mono'    : Monotone toFun
  modular' : ∀ a b, toFun a + toFun b
             = toFun (a ⊔ b) + toFun (a ⊓ b)
\end{lstlisting}

\noindent Proposition~\ref{prop:lvaluation-comap}.
\phantomsection\label{lean:lvaluation-comap}
\begin{lstlisting}
def LValuation.comap (f : BoundedLatticeHom Γ Ω) (v : LValuation Ω) :
    LValuation Γ

theorem LValuation.comap_id (v : LValuation Ω) :
    v.comap (BoundedLatticeHom.id Ω) = v

theorem LValuation.comap_comp (f : BoundedLatticeHom Γ Ω)
    (g : BoundedLatticeHom Δ Γ) (v : LValuation Ω) :
    v.comap (f.comp g) = (v.comap f).comap g
\end{lstlisting}

\noindent Theorem~\ref{thm:lvaluation-separating}.
\phantomsection\label{lean:lvaluation-separating}
\begin{lstlisting}
theorem LValuation.exists_sharp_separating [DistribLattice Ω]
    {a b : Ω} (hab : ¬ a ≤ b) :
    ∃ v : LValuation Ω, v.IsSharp ∧ v a = 1 ∧ v b = 0
\end{lstlisting}

\noindent Definition~\ref{def:m3}.
\phantomsection\label{lean:m3}
\begin{lstlisting}
inductive M3 : Type
  | z | a | b | c | o

instance : Lattice M3 where
  -- z bottom, o top, a b c pairwise-incomparable atoms;
  -- any two distinct atoms join to o and meet to z

instance : BoundedOrder M3 where
  top := o
  bot := z

theorem M3.not_distributive :
    ¬ ∀ x y w : M3, x ⊓ (y ⊔ w) = (x ⊓ y) ⊔ (x ⊓ w)
\end{lstlisting}

\noindent Theorem~\ref{thm:m3-no-sharp}.
\phantomsection\label{lean:m3-no-sharp}
\begin{lstlisting}
theorem M3.no_sharp_valuation (v : LValuation M3) : ¬ v.IsSharp
\end{lstlisting}

\noindent Definition~\ref{def:tensor-valuation}.
\phantomsection\label{lean:tensor-valuation}
\begin{lstlisting}
def IsTensorValuation (v : Ω → ℝ≥0∞) (t : Ω → Ω → Ω) : Prop :=
  ∀ a b, v a + v b = v (a ⊔ b) + v (t a b)
\end{lstlisting}

\noindent Theorem~\ref{thm:tensor-collapse}.
\phantomsection\label{lean:tensor-collapse}
\begin{lstlisting}
theorem IsTensorValuation.collapse {v : Ω → ℝ≥0∞} {t : Ω → Ω → Ω}
    (h : IsTensorValuation v t) {a : Ω} (ha : v a ≠ ⊤) :
    v (t a a) = v a
\end{lstlisting}

\noindent Example~\ref{ex:tensor-witness}.
\phantomsection\label{lean:tensor-witness}
\begin{lstlisting}
theorem forced_constant_of_add_valuation {v : ℕ → ℝ≥0∞}
    (h : IsTensorValuation v (· + ·)) (hfin : ∀ n, v n ≠ ⊤) (n : ℕ) :
    v n = v 0
\end{lstlisting}

\noindent Corollary~\ref{cor:tensor-dichotomy}.
\phantomsection\label{lean:tensor-dichotomy}
\begin{lstlisting}
theorem not_isTensorValuation_of_discriminates {v : Ω → ℝ≥0∞}
    {t : Ω → Ω → Ω} {a : Ω} (ha : v a ≠ ⊤) (hne : v (t a a) ≠ v a) :
    ¬ IsTensorValuation v t
\end{lstlisting}
\subsection{Probability Calculi from Certain States}

\phantomsection\label{lean:mix}
\begin{lstlisting}
def IsSharpVec (x : E → ℝ≥0∞) : Prop := ∀ e, x e = 0 ∨ x e = 1

def Mix (S : Set (E → ℝ≥0∞)) : Set (E → ℝ≥0∞) :=
  {x | ∃ (n : ℕ) (w : Fin n → ℝ≥0∞) (s : Fin n → E → ℝ≥0∞),
    (∀ i, s i ∈ S) ∧ ∑ i, w i = 1 ∧ x = ∑ i, w i • s i}
\end{lstlisting}

\phantomsection\label{lean:sharp-mem-mix-iff}
\begin{lstlisting}
theorem sharp_mem_Mix_iff {S : Set (E → ℝ≥0∞)} (hS : ∀ s ∈ S, IsSharpVec s)
    {x : E → ℝ≥0∞} (hx : IsSharpVec x) : x ∈ Mix S ↔ x ∈ S
\end{lstlisting}

\phantomsection\label{lean:mix-mix}
\begin{lstlisting}
theorem Mix_Mix (S : Set (E → ℝ≥0∞)) : Mix (Mix S) = Mix S
\end{lstlisting}

\phantomsection\label{lean:affinelaw}
\begin{lstlisting}


structure AffineLaw (E : Type*) [Fintype E] where
  a : E → ℝ≥0∞
  b : E → ℝ≥0∞
  k : ℝ≥0∞
  k' : ℝ≥0∞

def AffineLaw.Holds [Fintype E] (L : AffineLaw E) (x : E → ℝ≥0∞) : Prop :=
  ∑ e, L.a e * x e + L.k ≤ ∑ e, L.b e * x e + L.k'
\end{lstlisting}

\phantomsection\label{lean:affinelaw-holds-of-mix}
\begin{lstlisting}
theorem AffineLaw.holds_of_Mix [Fintype E] (L : AffineLaw E) {S : Set (E → ℝ≥0∞)}
    (hL : ∀ s ∈ S, L.Holds s) {x : E → ℝ≥0∞} (hx : x ∈ Mix S) : L.Holds x
\end{lstlisting}

\phantomsection\label{lean:mix-eq-laws}
\begin{lstlisting}
theorem Mix_eq_laws' {S : Set (E → ℝ≥0∞)} (hS : ∀ s ∈ S, IsSharpVec s) :
    Mix S = {x | ∀ L : AffineLaw E, (∀ s ∈ S, L.Holds s) → L.Holds x}
\end{lstlisting}

\phantomsection\label{lean:frame-recipe}
\begin{lstlisting}
theorem frame_recipe :
    Set.range (fun v : Valuation (LowerSet P) => v.toFun)
      = Mix {f | ∃ v : Valuation (LowerSet P), v.IsSharp ∧ f = v.toFun}
\end{lstlisting}

\phantomsection\label{lean:eq-deltapoint-of-issharp}
\begin{lstlisting}
theorem Valuation.eq_deltaPoint_of_isSharp {v : Valuation (LowerSet P)} (hv : v.IsSharp) :
    ∃ p, v = deltaPoint p
\end{lstlisting}

\phantomsection\label{lean:laws-iff-local-mix}
\begin{lstlisting}
structure FinLaw (E : Type*) where
  F : Finset E
  a : E → ℝ≥0∞
  b : E → ℝ≥0∞
  k : ℝ≥0∞
  k' : ℝ≥0∞

def FinLaw.Holds (L : FinLaw E) (x : E → ℝ≥0∞) : Prop :=
  ∑ e ∈ L.F, L.a e * x e + L.k ≤ ∑ e ∈ L.F, L.b e * x e + L.k'

def restr (F : Finset E) (x : E → ℝ≥0∞) : F → ℝ≥0∞ := fun e => x e

theorem laws_iff_local_Mix {S : Set (E → ℝ≥0∞)} (hS : ∀ s ∈ S, IsSharpVec s)
    (x : E → ℝ≥0∞) :
    (∀ L : FinLaw E, (∀ s ∈ S, L.Holds s) → L.Holds x) ↔
      ∀ F : Finset E, restr F x ∈ Mix (restr F '' S)
\end{lstlisting}

\phantomsection\label{lean:laws-iff-barycenter}
\begin{lstlisting}
def Cfg (Ev : Type) := Ev → Bool

def Good (S : Set (Ev → ℝ≥0∞)) (F : Finset Ev) (σ : Cfg Ev) : Prop :=
  ∃ s ∈ S, ∀ e : F, s e.1 = vec σ e.1

def atom (e : Ev) : Set (Cfg Ev) := {σ | σ e = true}

def bad (S : Set (Ev → ℝ≥0∞)) (F : Finset Ev) : Set (Cfg Ev) := {σ | ¬ Good S F σ}

theorem laws_iff_barycenter {S : Set (Ev → ℝ≥0∞)} (hS : ∀ s ∈ S, IsSharpVec s)
    (x : Ev → ℝ≥0∞) :
    (∀ L : FinLaw Ev, (∀ s ∈ S, L.Holds s) → L.Holds x) ↔
      ∃ μ : ProbabilityMeasure (Cfg Ev), (∀ F, (μ : Measure (Cfg Ev)) (bad S F) = 0) ∧
        ∀ e, x e = (μ : Measure (Cfg Ev)) (atom e)
\end{lstlisting}

\phantomsection\label{lean:mem-lin-iff}
\begin{lstlisting}
def flat (X : Type*) [Fintype X] : Set (X → ℝ≥0∞) := {x | ∑ a, x a ≤ 1}

noncomputable def app (t : X × Y → ℝ≥0∞) (x : X → ℝ≥0∞) : Y → ℝ≥0∞ :=
  fun b => ∑ a, x a * t (a, b)

def lin (X Y : Type*) [Fintype X] [Fintype Y] : Set (X × Y → ℝ≥0∞) :=
  {t | ∀ x ∈ flat X, app t x ∈ flat Y}

theorem mem_lin_iff [DecidableEq X] (t : X × Y → ℝ≥0∞) :
    t ∈ lin X Y ↔ ∀ a, ∑ b, t (a, b) ≤ 1
\end{lstlisting}

\phantomsection\label{lean:ispcs-lin}
\begin{lstlisting}
def orth (S : Set (E → ℝ≥0∞)) : Set (E → ℝ≥0∞) := {y | ∀ x ∈ S, pairing x y ≤ 1}

structure IsPCS (P : Set (E → ℝ≥0∞)) : Prop where
  closed : orth (orth P) = P
  total : ∀ e, ∃ x ∈ P, x e ≠ 0
  bounded : ∀ e, ∃ M ≠ ⊤, ∀ x ∈ P, x e ≤ M

theorem isPCS_lin : IsPCS (lin X Y)

theorem isPCS_flat : IsPCS (flat X)

theorem isPCS_tensDual : IsPCS tensDual
\end{lstlisting}

\phantomsection\label{lean:pcs-recipe}
\begin{lstlisting}
theorem pcs_recipe {X Y : Type} [Fintype X] [Fintype Y] [DecidableEq X] :
    lin X Y = Mix (Set.range (det : (X → Option Y) → X × Y → ℝ≥0∞))
\end{lstlisting}

\phantomsection\label{lean:sharp-lin-iff-clique}
\begin{lstlisting}
theorem sharp_lin_iff_clique {X Y : Type*} [Fintype X] [Fintype Y] [DecidableEq X] [DecidableEq Y]
    {t : X × Y → ℝ≥0∞} (h01 : ∀ p, t p = 0 ∨ t p = 1) :
    t ∈ lin X Y ↔ IsClique (arrowC (flatC X) (flatC Y)) (t · = 1)
\end{lstlisting}

\phantomsection\label{lean:sharp-td-iff}
\begin{lstlisting}
def tD (X Y X' Y' : Type) [Fintype X] [Fintype Y] [Fintype X'] [Fintype Y'] :
    Set ((X × Y) × (X' × Y') → ℝ≥0∞) :=
  {t | ∀ y ∈ lin X Y, ∀ z ∈ lin X' Y', ∑ p, ∑ q, y p * t (p, q) * z q ≤ 1}

def FTree (a : X) (cf : Y → X') (pq : (X × Y) × (X' × Y')) : Prop :=
  pq.1.1 = a ∧ pq.2.1 = cf pq.1.2

theorem sharp_tD_iff [Nonempty X] [Nonempty X'] {z : (X × Y) × (X' × Y') → ℝ≥0∞}
    (h01 : ∀ pq, z pq = 0 ∨ z pq = 1) :
    z ∈ tD X Y X' Y' ↔
      (∃ a cf, ∀ pq, z pq = 1 → FTree a cf pq) ∨ (∃ c af, ∀ pq, z pq = 1 → GTree c af pq)
\end{lstlisting}

\phantomsection\label{lean:sharp-mem-sequentialg}
\begin{lstlisting}
theorem sharp_mem_sequentialG [Nonempty X] [Nonempty X'] {z : (X × Y) × (X' × Y') → ℝ≥0∞}
    (hz : z ∈ tD X Y X' Y') (h01 : ∀ u, z u = 0 ∨ z u = 1) : z ∈ sequentialG X Y X' Y'
\end{lstlisting}

\phantomsection\label{lean:recipe-eq-sequentialg}
\begin{lstlisting}
structure Core (A B C D : Type) [Fintype A] [Fintype C] where
  μ : A → ℝ≥0∞
  κ : A → B → C → ℝ≥0∞
  ρ : A → B → C → D → ℝ≥0∞
  hκ : ∀ a b, ∑ c, κ a b c ≤ 1
  hρ : ∀ a b c d, ρ a b c d ≤ 1

def sequentialG (X Y X' Y' : Type) [Fintype X] [Fintype X'] : Set ((X × Y) × (X' × Y') → ℝ≥0∞) :=
  {t | ∃ (F : Core X Y X' Y') (G : Core X' Y' X Y), F.mass + G.mass ≤ 1 ∧
    t = fun pq => F.val pq.1.1 pq.1.2 pq.2.1 pq.2.2 + G.val pq.2.1 pq.2.2 pq.1.1 pq.1.2}

theorem recipe_eq_sequentialG [Nonempty X] [Nonempty X'] :
    Mix {z | z ∈ tD X Y X' Y' ∧ ∀ u, z u = 0 ∨ z u = 1} = sequentialG X Y X' Y'
\end{lstlisting}

\phantomsection\label{lean:sequentialg-subset-mix}
\begin{lstlisting}
theorem sequentialG_subset_Mix : sequentialG X Y X' Y' ⊆ Mix (sharpSeqG X Y X' Y')
\end{lstlisting}

\phantomsection\label{lean:mix-sequentialg}
\begin{lstlisting}
theorem Mix_sequentialG : Mix (sequentialG X Y X' Y') = sequentialG X Y X' Y'
\end{lstlisting}

\phantomsection\label{lean:pente-not-mix}
\begin{lstlisting}
theorem pentE_not_mix (hx : Function.Injective ex) (hy : Function.Injective ey)
    (hx' : Function.Injective ex') (hy' : Function.Injective ey')
    {n : ℕ} (w : Fin n → ℝ≥0∞) (hw : ∑ i, w i = 1) (z : Fin n → (X × Y) × (X' × Y') → ℝ≥0∞)
    (hz : ∀ i, z i ∈ tD X Y X' Y') (h01 : ∀ i u, z i u = 0 ∨ z i u = 1) :
    pentE ex ey ex' ey' ≠ ∑ i, w i • z i

theorem pentE_mem (hx : Function.Injective ex) (hy : Function.Injective ey)
    (hx' : Function.Injective ex') (hy' : Function.Injective ey') :
    pentE ex ey ex' ey' ∈ tD X Y X' Y'
\end{lstlisting}

\phantomsection\label{lean:pent-cycle}
\begin{lstlisting}
theorem pent_clique_le_two : ∀ B ∈ pent.powerset,
    (∀ s ∈ B, ∀ t ∈ B, compat2 s t) → B.card ≤ 2

theorem pent_indep_le_two : ∀ B ∈ pent.powerset,
    (∀ s ∈ B, ∀ t ∈ B, compat2 s t → s = t) → B.card ≤ 2
\end{lstlisting}

\phantomsection\label{lean:pent-bounds}
\begin{lstlisting}
def tensDual : Set (W × W → ℝ≥0∞) :=
  {t | ∀ y ∈ lin Bool Bool, ∀ z ∈ lin Bool Bool, ∑ p, ∑ q, y p * t (p, q) * z q ≤ 1}

def sequential : Set (W × W → ℝ≥0∞) :=
  {t | ∃ (α β : ℝ≥0∞) (F : FFirst) (G : GFirst), α + β ≤ 1 ∧ t = α • F.den + β • G.den}

theorem seq_pent_le_two {t : W × W → ℝ≥0∞} (ht : t ∈ sequential) : ∑ pq ∈ pent, t pq ≤ 2

theorem seq_pent_tight : ∃ t ∈ sequential, ∑ pq ∈ pent, t pq = 2

theorem pentElt_pent_sum : ∑ pq ∈ pent, pentElt pq = 5 * 2⁻¹

theorem tensDual_pent_le {t : W × W → ℝ≥0∞} (ht : t ∈ tensDual) :
    2 * ∑ pq ∈ pent, t pq ≤ 5
\end{lstlisting}

\phantomsection\label{lean:sharp-bang-mixed-zero}
\begin{lstlisting}
theorem sharp_bang_mixed_zero {z : ℕ × ℕ → ℝ≥0∞} (hz : z ∈ bang) (h01 : ∀ m, z m = 0 ∨ z m = 1)
    {i j : ℕ} (hi : 1 ≤ i) (hj : 1 ≤ j) : z (i, j) = 0
\end{lstlisting}

\phantomsection\label{lean:prom-half-not-mix}
\begin{lstlisting}
noncomputable def prom (x : Bool → ℝ≥0∞) : ℕ × ℕ → ℝ≥0∞ := fun m => x true ^ m.1 * x false ^ m.2

def bang : Set (ℕ × ℕ → ℝ≥0∞) := orthI (orthI gens)

theorem prom_half_not_mix : prom half ∈ bang ∧
    prom half ∉ Mix {z | z ∈ bang ∧ ∀ m, z m = 0 ∨ z m = 1}
\end{lstlisting}

\phantomsection\label{lean:td-eq-qstab}
\begin{lstlisting}
def excl (s t : (X × Y) × (X' × Y')) : Prop := s ≠ t ∧ cp2 s t

theorem tD_eq_QSTAB : tD X Y X' Y' = QSTAB (excl (X := X) (Y := Y) (X' := X') (Y' := Y'))
\end{lstlisting}

\phantomsection\label{lean:recipe-eq-stab}
\begin{lstlisting}
theorem recipe_eq_STAB :
    Mix {z | z ∈ tD X Y X' Y' ∧ ∀ u, z u = 0 ∨ z u = 1}
      = STAB (excl (X := X) (Y := Y) (X' := X') (Y' := Y'))
\end{lstlisting}

\phantomsection\label{lean:berge-of-stab-eq-qstab}
\begin{lstlisting}
def QSTAB : Set (V → ℝ≥0∞) := {x | ∀ C, IsCliqueG adj C → ∑ v ∈ C, x v ≤ 1}

def STAB : Set (V → ℝ≥0∞) := Mix {x | ∃ I, IsStableG adj I ∧ x = ind I}

structure Hole (adj : V → V → Prop) (n : ℕ) where
  h : ZMod n → V
  inj : Function.Injective h
  adj_iff : ∀ i j, i ≠ j → (adj (h i) (h j) ↔ cyc i j)

theorem berge_of_stab_eq_qstab (heq : STAB adj = QSTAB adj) (k : ℕ) (hk : 2 ≤ k)
    [NeZero (2 * k + 1)] : IsEmpty (Hole adj (2 * k + 1)) ∧ IsEmpty (Antihole adj (2 * k + 1))
\end{lstlisting}

\phantomsection\label{lean:berge-of-recipe-eq-pcs}
\begin{lstlisting}
theorem berge_of_recipe_eq_pcs
    (heq : Mix {z | z ∈ tD X Y X' Y' ∧ ∀ u, z u = 0 ∨ z u = 1} = tD X Y X' Y')
    (k : ℕ) (hk : 2 ≤ k) [NeZero (2 * k + 1)] :
    IsEmpty (Hole (excl (X := X) (Y := Y) (X' := X') (Y' := Y')) (2 * k + 1)) ∧
      IsEmpty (Antihole (excl (X := X) (Y := Y) (X' := X') (Y' := Y')) (2 * k + 1))
\end{lstlisting}

\subsection{Multiplicative--Additive Formulas}

\phantomsection\label{lean:linear-fm}
\begin{lstlisting}
inductive Fm where
  | base (n : ℕ)
  | neg (A : Fm)
  | tens (A B : Fm)
  | wth (A B : Fm)

def web : Fm → Type
  | base n => Fin n
  | neg A => web A
  | tens A B => web A × web B
  | wth A B => web A ⊕ web B

def P : (A : Fm) → Set (web A → ℝ≥0∞)
  | base n => flat (Fin n)
  | neg A => orth (P A)
  | tens A B => orth (orth {z | ∃ x ∈ P A, ∃ y ∈ P B, z = tvec x y})
  | wth A B => {z | (fun a => z (Sum.inl a)) ∈ P A ∧ (fun b => z (Sum.inr b)) ∈ P B}

def coh : (A : Fm) → web A → web A → Prop
  | base _ => fun x y => x = y
  | neg A => fun x y => x = y ∨ ¬ coh A x y
  | tens A B => fun s t => coh A s.1 t.1 ∧ coh B s.2 t.2
  | wth A B => fun s t => match s, t with
    | Sum.inl a, Sum.inl a' => coh A a a'
    | Sum.inr b, Sum.inr b' => coh B b b'
    | _, _ => True
\end{lstlisting}

\phantomsection\label{lean:ispcs-p}
\begin{lstlisting}
theorem isPCS_P : ∀ A : Fm, IsPCS (P A)
\end{lstlisting}

\phantomsection\label{lean:linear-invariants}
\begin{lstlisting}
theorem invariants : ∀ A : Fm,
    (∀ s : web A, Pi.single s 1 ∈ P A ∧ ∀ x ∈ P A, x s ≤ 1) ∧
    (∀ s t : web A, s ≠ t → coh A s t → pr s t ∈ P A) ∧
    (∀ s t : web A, s ≠ t → ¬ coh A s t → ∀ x ∈ P A, x s + x t ≤ 1)

theorem single_mem_P (A : Fm) (s : web A) : Pi.single s 1 ∈ P A

theorem pair_mem_P (A : Fm) {s t : web A} (hst : s ≠ t) (h : coh A s t) : pr s t ∈ P A
\end{lstlisting}

\phantomsection\label{lean:sharp-isclique}
\begin{lstlisting}
theorem sharp_isClique (A : Fm) {z : web A → ℝ≥0∞} (hz : z ∈ P A)
    (h01 : ∀ s, z s = 0 ∨ z s = 1) : ∀ s t, z s = 1 → z t = 1 → coh A s t
\end{lstlisting}

\phantomsection\label{lean:recipe-subset-p}
\begin{lstlisting}
theorem recipe_subset_P (A : Fm) :
    Mix {z | z ∈ P A ∧ ∀ s, z s = 0 ∨ z s = 1} ⊆ P A

theorem recipe_sharp_iff (A : Fm) {x : web A → ℝ≥0∞} (hx : ∀ s, x s = 0 ∨ x s = 1) :
    x ∈ Mix {z | z ∈ P A ∧ ∀ s, z s = 0 ∨ z s = 1} ↔ x ∈ P A
\end{lstlisting}

\phantomsection\label{lean:mix-subset-orth}
\begin{lstlisting}
theorem Mix_subset_orth {S T : Set (E → ℝ≥0∞)} (h : S ⊆ orth T) : Mix S ⊆ orth T
\end{lstlisting}

\phantomsection\label{lean:p-lolli-base}
\begin{lstlisting}
theorem P_lolli_base (n m : ℕ) : P (lolli (base n) (base m)) = lin (Fin n) (Fin m)
\end{lstlisting}

\phantomsection\label{lean:p-second}
\begin{lstlisting}
def second (n m n' m' : ℕ) : Fm := neg (tens (lolli (base n) (base m)) (lolli (base n') (base m')))

theorem P_second (n m n' m' : ℕ) :
    P (second n m n' m') = tD (Fin n) (Fin m) (Fin n') (Fin m')
\end{lstlisting}

\phantomsection\label{lean:recipe-formula}
\begin{lstlisting}
theorem recipe_first_order (n m : ℕ) :
    P (lolli (base n) (base m)) = ConstructiveProb.Mix (Set.range
      (det : (Fin n → Option (Fin m)) → Fin n × Fin m → ℝ≥0∞))

theorem recipe_second_order {n m n' m' : ℕ} (hn : 0 < n) (hn' : 0 < n') :
    ConstructiveProb.Mix {z | z ∈ P (second n m n' m') ∧ ∀ u, z u = 0 ∨ z u = 1}
      = sequentialG (Fin n) (Fin m) (Fin n') (Fin m')

theorem recipe_second_order_ssubset {n m n' m' : ℕ} (hn : 2 ≤ n) (hm : 2 ≤ m) (hn' : 2 ≤ n')
    (hm' : 2 ≤ m') :
    ConstructiveProb.Mix {z | z ∈ P (second n m n' m') ∧ ∀ u, z u = 0 ∨ z u = 1}
      ⊂ P (second n m n' m')
\end{lstlisting}

\phantomsection\label{lean:shannon-clique}
\begin{lstlisting}
abbrev S2 : Fm := second 2 2 2 2

abbrev TT : Fm := neg (tens S2 S2)

noncomputable def shannonF : Finset (web (tens S2 S2)) := shannon.image fun u => (φ u.1, φ u.2)

theorem shannonF_clique : ∀ u ∈ shannonF, ∀ v ∈ shannonF, coh TT u v
\end{lstlisting}

\phantomsection\label{lean:shannon-not-mem}
\begin{lstlisting}
theorem shannonF_not_mem : (fun u => if u ∈ shannonF then (1 : ℝ≥0∞) else 0) ∉ P TT
\end{lstlisting}

\subsection{Quantum Event Logic}

\phantomsection\label{lean:cabello-set}
\begin{lstlisting}
def vec : Fin 18 → Fin 4 → ℤ := ![
  ![-1, 1, 1, 1], ![0, 0, 0, 1], ![0, 0, 1, 0], ![0, 0, 1, 1], ![0, 1, -1, 0], ![0, 1, 0, -1],
  ![0, 1, 0, 0], ![1, -1, -1, 1], ![1, -1, 0, 0], ![1, -1, 1, -1], ![1, 0, -1, 0], ![1, 0, 0, -1],
  ![1, 0, 0, 1], ![1, 0, 1, 0], ![1, 1, -1, 1], ![1, 1, 0, 0], ![1, 1, 1, -1], ![1, 1, 1, 1]]

def bas : Fin 9 → Fin 4 → Fin 18 := ![
  ![1, 2, 15, 8], ![1, 6, 13, 10], ![9, 7, 15, 3], ![9, 17, 10, 5], ![2, 6, 12, 11],
  ![7, 17, 11, 4], ![14, 16, 8, 3], ![14, 0, 13, 5], ![16, 0, 12, 4]]

def dot (u v : Fin 4 → ℤ) : ℤ := ∑ k, u k * v k
\end{lstlisting}

\phantomsection\label{lean:basis-orth}
\begin{lstlisting}
theorem basis_orth : ∀ b : Fin 9, ∀ i j : Fin 4, i ≠ j → dot (vec (bas b i)) (vec (bas b j)) = 0

theorem occ_two : ∀ v : Fin 18,
    (univ.filter (fun p : Fin 9 × Fin 4 => bas p.1 p.2 = v)).card = 2
\end{lstlisting}

\phantomsection\label{lean:no-ks}
\begin{lstlisting}
theorem no_ks : ¬ ∃ f : Fin 18 → Bool, ∀ b, (univ.filter (fun i => f (bas b i))).card = 1
\end{lstlisting}

\phantomsection\label{lean:no-certain-state}
\begin{lstlisting}
def certain : Set (Fin 18 → ℝ≥0∞) :=
  {s | IsSharpVec s ∧ ∀ b, ∑ i, s (bas b i) = 1}

theorem no_certain_state : certain = ∅

theorem Mix_certain_empty : Mix certain = ∅
\end{lstlisting}

\phantomsection\label{lean:parseval}
\begin{lstlisting}
theorem parseval (ψ : Fin 4 → ℝ) (b : Fin 9) :
    ∑ i, (∑ k, ψ k * vec (bas b i) k) ^ 2 / ∑ k, ((vec (bas b i) k : ℝ)) ^ 2 = ∑ k, ψ k ^ 2
\end{lstlisting}

\phantomsection\label{lean:born-not-mix}
\begin{lstlisting}
def states : Set (Fin 18 → ℝ≥0∞) := {x | ∀ b, ∑ i, x (bas b i) = 1}

noncomputable def born (ψ : Fin 4 → ℝ) : Fin 18 → ℝ≥0∞ := fun v =>
  ENNReal.ofReal ((∑ k, ψ k * vec v k) ^ 2 / ∑ k, ((vec v k : ℝ)) ^ 2)

theorem born_mem {ψ : Fin 4 → ℝ} (hψ : ∑ k, ψ k ^ 2 = 1) : born ψ ∈ states

theorem born_not_mix {ψ : Fin 4 → ℝ} (hψ : ∑ k, ψ k ^ 2 = 1) :
    born ψ ∈ states ∧ born ψ ∉ Mix certain
\end{lstlisting}

\subsection{Finite Frames without Choice}

\phantomsection\label{lean:exists-sharp-separating-fin}
\begin{lstlisting}
def JoinPrime (j : α) : Prop := ∀ x y : α, j ≤ x ⊔ y → j ≤ x ∨ j ≤ y

theorem exists_joinPrime_aux : ∀ (n : ℕ) (a b : α), (below a).card = n → ¬ a ≤ b →
    ∃ j, j ≤ a ∧ ¬ j ≤ b ∧ JoinPrime j

theorem exists_sharp_separating_fin [BoundedOrder α] {a b : α} (hab : ¬ a ≤ b) :
    ∃ v : α → ℕ, (∀ x, v x = 0 ∨ v x = 1) ∧ Monotone v ∧ v ⊥ = 0 ∧ v ⊤ = 1 ∧
      (∀ x y, v x + v y = v (x ⊔ y) + v (x ⊓ y)) ∧ v a = 1 ∧ v b = 0
\end{lstlisting}

\phantomsection\label{lean:em-of-sharp-complement-fin}
\begin{lstlisting}
theorem em_of_sharp_complement_fin [BoundedOrder α] (a c : α)
    (h : ∀ v : α → ℕ, (∀ x, v x = 0 ∨ v x = 1) → Monotone v → v ⊥ = 0 → v ⊤ = 1 →
      (∀ x y, v x + v y = v (x ⊔ y) + v (x ⊓ y)) → v a + v c = 1) :
    a ⊔ c = ⊤
\end{lstlisting}

\subsection{Choice-Free Finite Results}

\phantomsection\label{lean:constr-sierpinski}
\begin{lstlisting}
inductive Op | bot | halts | top
  deriving DecidableEq

def neg : Op → Op
  | bot => top
  | halts => bot
  | top => bot

theorem neg_spec : ∀ o o' : Op, (inf o' o = bot) ↔ le o' (neg o) = true

theorem neg_halts : neg halts = bot

theorem frontier_iff : ∀ x : Bool, frontier x = true ↔ x = false

theorem slack_eq (wT wF : Nat) :
    val wT wF top = val wT wF halts + val wT wF (neg halts) + wt (frontier false) wF
\end{lstlisting}

\phantomsection\label{lean:constr-boundary}
\begin{lstlisting}
def ext (ops : List (Fin n → Bool)) (U : Fin n → Bool) (x : Fin n) : Bool :=
  ops.any fun O => O x && allFin n (fun y => !(O y && U y))

def bd (ops : List (Fin n → Bool)) (U : Fin n → Bool) (x : Fin n) : Bool :=
  !U x && !ext ops U x

theorem sub_ext {ops : List (Fin n → Bool)} {U : Fin n → Bool} {O : Fin n → Bool}
    (hO : O ∈ ops) (hd : ∀ y, (!(O y && U y)) = true) {x : Fin n} (hx : O x = true) :
    ext ops U x = true

theorem slack_eq_boundary_fin (ops : List (Fin n → Bool)) (U : Fin n → Bool) (μ : Fin n → Nat) :
    fs n μ = fs n (fun x => wt (U x) (μ x)) + fs n (fun x => wt (ext ops U x) (μ x)) +
      fs n (fun x => wt (bd ops U x) (μ x))
\end{lstlisting}

\phantomsection\label{lean:constr-kraft}
\begin{lstlisting}
def ksum (L : Nat) : List (List Bool) → Nat
  | [] => 0
  | w :: S => 2 ^ (L - w.length) + ksum L S

theorem kraft : ∀ (L : Nat) (S : List (List Bool)), PF S → (∀ w ∈ S, w.length ≤ L) →
    ksum L S ≤ 2 ^ L
\end{lstlisting}

\phantomsection\label{lean:constr-farkas}
\begin{lstlisting}
structure Row (n : Nat) where
  c : Fin n → Int
  s : Bool

inductive Der {n : Nat} (rows : List (Row n)) : Row n → Prop
  | base {r : Row n} : r ∈ rows → Der rows r
  | comb {r₁ r₂ : Row n} (p q : Int) : 0 ≤ p → 0 ≤ q → Der rows r₁ → Der rows r₂ →
      Der rows (comb p q r₁ r₂)

theorem fm : ∀ (n : Nat) (rows : List (Row n)), (∃ w, ∀ r ∈ rows, Sat r w) ∨ Der rows (zeroS n)

theorem complete (hd : 0 < d)
    (hlaw : ∀ (c : Fin m → Int) (k : Int), (∀ j, k ≤ dot m c (S j)) → k * d ≤ dot m c y) :
    ∃ (lam : Fin N → Int) (t : Int), 0 < t ∧ (∀ j, 0 ≤ lam j) ∧
      dot N (fun _ => 1) lam = t ∧ ∀ e, d * dot N lam (fun j => S j e) = t * y e

theorem sound (hd : 0 < d) (lam : Fin N → Int) (t : Int) (ht : 0 < t) (hl : ∀ j, 0 ≤ lam j)
    (hsum : dot N (fun _ => 1) lam = t) (hx : ∀ e, d * dot N lam (fun j => S j e) = t * y e)
    (c : Fin m → Int) (k : Int) (hc : ∀ j, k ≤ dot m c (S j)) : k * d ≤ dot m c y
\end{lstlisting}

\phantomsection\label{lean:constr-fubini}
\begin{lstlisting}
structure EOD (A : Nat → Nat → Bool) where
  N : Nat
  lt : Bool
  eq : Bool
  gt : Bool
  spec : ∀ m m', N ≤ m → N ≤ m' →
    A m m' = if m < m' then lt else if m = m' then eq else gt

def L {A : Nat → Nat → Bool} (w : EOD A) : Bool := w.lt
/-- The iterated limit with `m → ∞` first. -/
def R {A : Nat → Nat → Bool} (w : EOD A) : Bool := w.gt

def R {A : Nat → Nat → Bool} (w : EOD A) : Bool := w.gt

theorem L_modular {A B} (a : EOD A) (b : EOD B) :
    val (L a) + val (L b) = val (L (union a b)) + val (L (inter a b))

theorem rect_L {S T} (s : EC S) (t : EC T) : L (rect s t) = (s.v && t.v)

theorem rect_R {S T} (s : EC S) (t : EC T) : R (rect s t) = (s.v && t.v)

theorem products_differ : L tri ≠ R tri
\end{lstlisting}

\section{The Acz\'el Development}
\label{app:aczel}

This appendix sketches the from-scratch construction behind the
product-rule half of Theorem~\ref{thm:cox}.
The target shape is Acz\'el's associativity
theorem~\cite{aczel1966}: a suitably regular associative
$F : \mathbb{R} \to \mathbb{R} \to \mathbb{R}$ regraduates to
multiplication, $G(F\,x\,y) = G(x) \cdot G(y)$ for all $x, y$ and a
strictly monotone $G$.  \textsf{mathlib} has no
associativity representation, no H\"older-style embedding of ordered
semigroups on intervals, and no Cauchy functional equation under
monotonicity (only under continuity).

\paragraph{Hypotheses.}
A \texttt{Scale} bundles: associativity of $F$, \emph{divisibility}
$\forall c\, \exists x,\ F\,x\,x = c$, positivity on the cone
$x < F\,x\,c$, joint order-preservation, joint continuity, a unit
$u$, and identity at $-\infty$, meaning
$F\,x\,z \to x$ as $z \to -\infty$.  Order-preservation is required
by the additivity proof below.  Divisibility yields the
$2^{n}$-th roots used in the dyadic construction below.  Without it
the generator is not definable.  It is assumed directly in the
bundle, though its interval form follows from the intermediate value
theorem (\texttt{exists\_diag\_eq}).  Non-vacuity: \emph{log-sum-exp},
$F\,x\,y = \log(e^x + e^y)$ with $u = 0$, satisfies every field of
the bundle (\texttt{logSumExpScale}), with intended generator
globally $\exp$.

\paragraph{Construction.}
Write $a^{(n)}$ for $n$ right-associated copies of $a$ under $F$
(\texttt{Fpow}).  Associativity alone makes
$n \mapsto a^{(n)}$ a semigroup morphism from $(\mathbb{N}, +)$.
Commutativity is not assumed, and follows instead as a corollary of
the generator.  Unboundedness of the iterates is derived: 
were they bounded, monotone convergence would give a fixed
point $F\,L\,a = L$, contradicting positivity, and unboundedness
yields the Archimedean property.  Divisibility yields a tower of
$2^n$-th roots $u \succ u^{1/2} \succ u^{1/4} \succ \cdots$
(\texttt{droot}), where $2^\ell$ copies of the finer root compose
under $F$ to one copy of the coarser root.  Counting the number of
copies of the $n$-th root not exceeding $b$
(\texttt{dcount}, well defined by the Archimedean property) gives
dyadic approximants $g_n(b) = \mathrm{dcount}(b,n)/2^n$.  The
refinement inequality $2N + 1 \le N' \le 2N + 2$, where $N$ and $N'$
are the counts at consecutive refinement levels $n$ and $n+1$, makes
$(g_n)$ monotone and Cauchy, and the generator is its limit,
$g(b) = \bigsqcup_n g_n(b)$.

\paragraph{Results (all \texttt{sorry}-free, standard axioms).}
On the cone $\{x \mid u \le x\}$ we prove three things.
Normalization: $g(u) = 1$.  Additivity:
$g(F\,x\,y) = g(x) + g(y)$, since
$\mathrm{dcount}\,x + \mathrm{dcount}\,y + 1 \le
\mathrm{dcount}(F\,x\,y) \le \mathrm{dcount}\,x +
\mathrm{dcount}\,y + 3$ and the $2^{-n}$ error term vanishes.  And
strict monotonicity, proved without any continuity of $g$: a root
inserted between $x < y$ contributes a full extra block of
$2^{n-M}$ copies, $M$ being the refinement level of the inserted
root.  Exponentiating gives the multiplicative form, a strictly
monotone $G$ with $G(F\,x\,y) = G(x)\,G(y)$ on the cone
(\texttt{exists\_mul\_generator}), the conclusion of Acz\'el's
theorem restricted to the cone
(\texttt{aczelStatement\_cone}).  The reorientation
$\bar G = \exp(-g)$ gives the $[0,1]$-valued form used in
Section~\ref{sec:cox}.

\paragraph{Remaining gap.}
The constructed $g$ is discontinuous at the unit (it is $0$ below
$u$), so matching the classical statement verbatim needs the
off-cone extension (a group completion) and continuity of the
extension.  This is open.  It is not vacuous: the log-sum-exp
example admits a global continuous generator
(\texttt{hasOrderedGenerator\_\allowbreak logSumExp}).  Of H\"older's
classical hypotheses, the Archimedean property is proved, 
and is what enables the embedding into $\mathbb{R}$.
The interval form of divisibility follows from the intermediate
value theorem.  Only the embedding's global regularity remains.


\end{document}
